\documentclass[11pt]{article}

\usepackage[margin=1in]{geometry}
\usepackage{amsmath, amssymb, amsthm}
\usepackage{booktabs}
\usepackage{longtable}
\usepackage{hyperref}
\usepackage{enumitem}
\usepackage{xcolor}
\usepackage{array}
\newcolumntype{P}[1]{>{\raggedright\arraybackslash}p{#1}}
\IfFileExists{microtype.sty}{\usepackage[protrusion=true,expansion=true]{microtype}}{}
\setlength{\emergencystretch}{3em}
\hbadness=10000
\hfuzz=2pt

% Common symbols
\newcommand{\kQ}{\kappa_{Q}}
\newcommand{\Ch}{C_{h}}
\newcommand{\Sh}{\Sigma_{h}}
\newcommand{\E}{\mathbb{E}}
\newcommand{\Rb}{\mathbb{R}}
\newcommand{\Nb}{\mathbb{N}}
\newcommand{\Zb}{\mathbb{Z}}
\newcommand{\Nrm}{\mathcal{N}}
\newcommand{\Studt}{\mathrm{t}}
\newcommand{\Hcal}{\mathcal{H}}
\newcommand{\Pcal}{\mathcal{P}}
\newcommand{\Mcal}{\mathcal{M}}
\newcommand{\Lcal}{\mathcal{L}}
\newcommand{\Bcal}{\mathcal{B}}
\newcommand{\PIT}{\mathrm{PIT}}
\newcommand{\KL}{\mathrm{KL}}
\newcommand{\defect}{\Delta}
\newcommand{\compose}{\star}
\DeclareMathOperator{\Var}{Var}
\DeclareMathOperator{\Cov}{Cov}

\theoremstyle{definition}
\newtheorem{definition}{Definition}
\newtheorem{remark}{Remark}
\theoremstyle{plain}
\newtheorem{proposition}{Proposition}
\newtheorem*{hypothesis}{Working hypothesis}

\title{\textbf{Relaxed Kernel Factorization:\\
A Kernel-Family-Agnostic Account of the Missing-Content Programme}\\[4pt]
\large A note: the structural framework, its Gaussian special case
as worked example, and what is settled at each level of generality}
\author{MITACS dynamics project --- internal memo, not a result artifact}
\date{\today}

\begin{document}

\maketitle

\begin{abstract}
\noindent
The seed \emph{Multi-Scale Factor Coherence via the Residual
Diffusion} originally proposed a relaxed-semigroup framework
for the $\kQ$-estimator with the diffusion identity as the
cross-scale coherence constraint. That formulation was implicitly
Gaussian. The empirical campaign succeeded in confirming the
framework's central claim on Ontario data under the Gaussian case
($\mathrm{D\_RATIO} = 28.96$, $95\%$ CI $[26.82, 31.43]$ against
the synthetic-VAR(1) noise floor) but also surfaced that the
Gaussian forecast distribution is itself operationally wrong
($\chi^{2}_{10\text{-bin}} = 22{,}329$). This note re-states the
framework at the right level of generality --- kernel-family-agnostic
--- so that the Gaussian case is a worked example and the non-Gaussian
extension is a natural object to study. The abstract framework, the
Gaussian special case, and the non-Gaussian extension all share the
same algebraic statement; only the test statistics change with the
kernel family. Q1 of the distributional-class thread has now
empirically tested the smallest non-trivial extension (Student-$t$):
the kernel-family lift reduces the marginal $\chi^{2}$ by $16\times$
($22{,}328 \to 1{,}371$) but does not fully calibrate, and the
cross-scale coherence defect at the v2 headline cell shrinks by
$\sim 10\times$ to $\mathrm{D\_RATIO}_{\Studt} = 2.93 \in [2.62,
3.24]$. Picture (C) survives the family lift at substantially
reduced magnitude; v2's settled finding is retained with the
asterisk that the magnitude was kernel-family-conditional. The
heavy-tailedness lives in the $\hat\kappa_1$ innovation, not the
seasonal climatology; Student-$t$ alone is not the full predictive
class, motivating richer-family experiments under the now-active
thread.
\end{abstract}

\section{Setting: process, embedding, kernel family}
\label{sec:setting}

Let $(z_t)_{t\in\Zb}$ be a real-valued stochastic process observed
at integer times, and let $x_t \in \Rb^d$ denote its $d$-dimensional
delay-coordinate embedding $x_t = (z_{t-1}, z_{t-2}, \ldots, z_{t-d})^{\top}$.
We make no parametric assumption about the law of $(z_t)$ a priori.

Let $\Bcal(\Rb^d)$ denote the Borel sigma-algebra on $\Rb^d$ and
$\Pcal(\Rb^d)$ the space of probability laws thereon. A
\emph{conditional Markov kernel} from $\Rb^d$ to itself is a
measurable map
\[
  \kappa : \Rb^d \to \Pcal(\Rb^d),
  \qquad x \mapsto \kappa(x, \cdot).
\]
The \emph{true} one-step conditional law of the process under the
embedding is the kernel $\kappa^{(1)}$ such that
$x_{t+1} \mid x_t \sim \kappa^{(1)}(x_t, \cdot)$, and the
$h$-step kernel is $\kappa^{(h)}$ for $h \in \Nb$.

If the process is Markov in the embedding, the $h$-step kernel is
determined by Chapman--Kolmogorov:
\begin{equation}
  \label{eq:CK}
  \kappa^{(h_1 + h_2)}(x, A) \;=\;
  \int_{\Rb^d} \kappa^{(h_2)}(y, A)\, \kappa^{(h_1)}(x, dy)
  \qquad \forall x \in \Rb^d,\ A \in \Bcal(\Rb^d).
\end{equation}
Define the kernel convolution $\compose$ by
$(\kappa_a \compose \kappa_b)(x, A) := \int \kappa_b(y, A)\,
\kappa_a(x, dy)$. Then \eqref{eq:CK} reads
$\kappa^{(h_1 + h_2)} = \kappa^{(h_1)} \compose \kappa^{(h_2)}$.
\emph{This statement is measure-theoretic, not parametric}: it
makes no commitment to a kernel family.

\section{The relaxed-factorization framework}
\label{sec:framework}

In practice, the true kernel $\kappa^{(h)}$ is unknown; we
estimate empirical kernels $\hat\kappa_h$ from finite data for
$h \in \Hcal \subseteq \Nb$. We consider three a-priori pictures of
how the family $\{\hat\kappa_h\}_{h \in \Hcal}$ might be structured.

\begin{definition}[strict / trivial / relaxed factorization]
\label{def:factorization}
The family $\{\hat\kappa_h\}_{h \in \Hcal}$ is:
\begin{description}[leftmargin=2em]
\item[(A) Strictly factored] if Chapman--Kolmogorov holds
exactly on the estimated kernels:
$\hat\kappa_{h_1 + h_2} = \hat\kappa_{h_1} \compose \hat\kappa_{h_2}$
for all $h_1, h_2 \in \Hcal$ with $h_1 + h_2 \in \Hcal$.

\item[(B) Trivially factored] (or \emph{unfactored}) if the
$\hat\kappa_h$ are unrelated --- each is estimated independently
and there is no algebraic relation imposed across scales.

\item[(C) Relaxedly factored] if there exists a family of
signed-measure-valued \emph{defect maps}
$\defect_{h_1, h_2} : \Rb^d \to \Mcal(\Rb^d)$, where $\Mcal(\Rb^d)$
denotes signed measures on $\Bcal(\Rb^d)$, such that
\begin{equation}
  \label{eq:relax}
  \hat\kappa_{h_1 + h_2}(x, \cdot) \;=\;
  (\hat\kappa_{h_1} \compose \hat\kappa_{h_2})(x, \cdot)
  \;+\; \defect_{h_1, h_2}(x, \cdot)
\end{equation}
for all $h_1, h_2 \in \Hcal$ with $h_1 + h_2 \in \Hcal$. The
strict factorization is the special case $\defect_{h_1, h_2} \equiv 0$;
the trivial factorization is the case where $\defect$ is
unconstrained.
\end{description}
\end{definition}

\begin{remark}[the role of the defect]
The defect $\defect_{h_1, h_2}$ is a signed measure (not a
probability) because $\hat\kappa_{h_1+h_2}$ and
$\hat\kappa_{h_1} \compose \hat\kappa_{h_2}$ are both probability
measures, and the difference is generally signed with total mass
zero. The defect is the natural carrier of all information about
how the family's empirical structure departs from the strict
semigroup. The seed's central claim is that for a finite-dimensional
embedding of a higher-dimensional process, the defect is non-trivial
\emph{and} carries identifiable structure about the dimensions of
variation the embedding does not span.
\end{remark}

\begin{definition}[coherence content]
\label{def:coherence}
The \emph{coherence content} of a relaxed factorization is the
structural information in $\{\defect_{h_1, h_2}\}$. Specifically:
any decomposition of the defect that is consistent across
$(h_1, h_2)$ pairs --- by origin-state dependence, by eigenmode
projection of an underlying drift, by horizon scale, by exogenous
stratification --- constitutes a partial characterization of the
embedding's missing degrees of freedom.
\end{definition}

\begin{hypothesis}[the seed's working claim]
\label{hyp:working}
For an empirical process whose true dynamics involve $k$ degrees
of freedom and whose embedding spans $d \leq k$ of them, the
relaxed factorization is the natural framework for inference: the
defects $\defect_{h_1, h_2}$ encode information about the $k - d$
unmodelled degrees of freedom, including their characteristic
timescales and their coupling to the embedded state.
\end{hypothesis}

This hypothesis is the empirical-mechanics conjecture behind the
seed; it can be supported or falsified by measuring the structure
of $\defect$. It is itself agnostic to kernel family.

\subsection{The structural role of the drift versus the rest of the kernel}
\label{sec:drift-vs-rest}

Even though $\hat\kappa_h$ is a full conditional law (not just a
mean), it is useful to decompose the strict-semigroup constraint
\eqref{eq:CK} into the \emph{drift identity} (constraint on
conditional means) and the \emph{everything else} (constraint on
the conditional law given the means).

The drift identity. If we write $\bar\mu_h(x) := \E[X' \mid X = x]$
under $\hat\kappa_h$, then strict factorization implies
$\bar\mu_{h_1 + h_2}(x) = \bar\mu_{h_2}\!\bigl(\bar\mu_{h_1}(x)\bigr)$.
For linear drifts $\bar\mu_h(x) = \hat C_h x$, this reduces to
$\hat C_{h_1 + h_2} = \hat C_{h_1} \hat C_{h_2}$, i.e.\ the drifts
form a semigroup under matrix multiplication. \emph{This is
automatic for any drift family that is closed under composition.}
A free-parameter linear drift is automatically closed; so are most
parametric drift families of interest. The drift identity therefore
imposes a weak constraint --- ``the estimated drifts compose as
matrix powers of \emph{something}'' --- and is uninformative on its
own about whether the strict factorization holds.

The everything-else identity. The non-drift information in $\hat\kappa_h$
--- its dispersion, its shape, its tail behaviour --- is where the
strict semigroup actually has empirical content. For each kernel
family this yields a different identity:

\begin{itemize}[leftmargin=2em]
\item For Gaussian kernels, the strict semigroup reduces to a
second-moment identity \eqref{eq:gauss-sigma} below.
\item For Student-$t$ kernels, the family is not closed under
convolution; the strict semigroup imposes that
$\hat\kappa_{h_1} \compose \hat\kappa_{h_2}$ (computed via numerical
convolution or sampling) equals $\hat\kappa_{h_1+h_2}$ as
distributions, not just as moments.
\item For non-parametric kernels (sample-based predictives), the
strict semigroup is a distributional equality between two
sample-based laws, tested via sample-based divergences
(Wasserstein, energy distance, MMD).
\end{itemize}

In all cases, \textbf{the structural role of the constraint is the
same}: the everything-else identity is the non-trivial part of the
strict factorization, and the defect $\defect_{h_1, h_2}$ is
operationalised by comparing the empirical
$\hat\kappa_{h_1 + h_2}$ to the convolved $\hat\kappa_{h_1} \compose
\hat\kappa_{h_2}$ under whichever divergence is appropriate for the
family.

\subsection{The forward equation as the underlying structural picture}
\label{sec:forward}

A continuous-time perspective makes the relationship between strict
factorization and an underlying generator transparent. For a
continuous-time Markov process with infinitesimal generator
$\Lcal$ (in the Markov-semigroup sense, acting on bounded measurable
functions), the family of $h$-step kernels satisfies the
\emph{Kolmogorov forward equation}
\begin{equation}
  \label{eq:forward}
  \partial_h \kappa^{(h)}(x, \cdot) \;=\; \Lcal^{*} \kappa^{(h)}(x, \cdot)
\end{equation}
where $\Lcal^{*}$ is the formal adjoint of $\Lcal$ (the
\emph{forward generator}, acting on measures). Strict factorization
of a discrete-time family $\{\hat\kappa_h\}$ corresponds, at the
ideal-infinite-data limit, to the existence of a single $\Lcal$
such that $\hat\kappa^{(h)} = e^{h \Lcal}$ (under appropriate
operator topology) consistently for all $h$.

We do not need to estimate $\Lcal$ to test strict factorization;
the test in discrete time is precisely the convolution test
\eqref{eq:relax}. The forward-equation framing is the theoretical
picture that explains \emph{why} the convolution test is the right
test, and what kind of object the defect $\defect$ is when strict
factorization fails (a non-generator-consistent residual in the
$h$-flow of the family of empirical kernels).

For the Gaussian case, $\Lcal^{*}$ has the form
$\Lcal^{*} f = -\nabla \cdot (b\, f) + \tfrac{1}{2}
\sum_{ij} \partial_i \partial_j(Q_{ij}\, f)$
(diffusion with drift $b$ and infinitesimal covariance $Q$).
For non-Gaussian Markov processes with jumps, $\Lcal^{*}$ includes
non-local integral operators (the L\'evy kernel). The framework's
structural content does not depend on which form $\Lcal^{*}$ takes;
only the closed-form expression for $\hat\kappa^{(h)}$ does.

\section{The Gaussian case as worked example}
\label{sec:gaussian}

We now specialize to the Gaussian kernel family. This is the case
the current registered estimator uses; the experiments described
in \S\ref{sec:settled} test the framework within this case.

\subsection{Gaussian kernel family and the strict semigroup}
\label{sec:gauss-strict}

Suppose $\hat\kappa_h$ is in the Gaussian family with linear drift:
\begin{equation}
  \label{eq:gauss-kernel}
  \hat\kappa_h(x, \cdot) \;=\; \Nrm\!\bigl(\hat C_h x,\; \hat\Sigma_h\bigr)
\end{equation}
for some $\hat C_h \in \Rb^{d \times d}$ and
$\hat\Sigma_h \in \Rb^{d \times d}$ (positive semidefinite). The
Gaussian family is closed under convolution: a linear-mean Gaussian
convolved with a linear-mean Gaussian is a linear-mean Gaussian.
Direct computation gives
\begin{equation}
  \label{eq:gauss-compose}
  (\hat\kappa_a \compose \hat\kappa_b)(x, \cdot) \;=\;
  \Nrm\!\bigl(\hat C_b \hat C_a x,\; \hat C_b \hat\Sigma_a \hat C_b^{\top} + \hat\Sigma_b\bigr).
\end{equation}
Iterating one-step kernels gives the closed-form parameters
\begin{align}
  \label{eq:gauss-Ch}
  \hat C_h^{\mathrm{iter}} &= \hat C_1^{\,h}, \\
  \label{eq:gauss-sigma}
  \hat\Sigma_h^{\mathrm{iter}} &= \sum_{j=0}^{h-1}
    \hat C_1^{\,j}\, \hat\Sigma_1\, (\hat C_1^{\,j})^{\top}.
\end{align}
Identity \eqref{eq:gauss-sigma} is the discrete-time analogue of
the continuous Lyapunov-type integral
\begin{equation}
  \label{eq:lyapunov}
  \Sigma_h \;=\; \int_0^h e^{sL}\, Q\, e^{sL^{\top}}\, ds
\end{equation}
that an Ornstein--Uhlenbeck SDE with drift $L$ and infinitesimal
covariance $Q$ would induce. We call \eqref{eq:gauss-sigma} (or
equivalently \eqref{eq:lyapunov}) the \emph{Gaussian diffusion
identity}. In the Gaussian case it is the everything-else
constraint of \S\ref{sec:drift-vs-rest}.

\subsection{Gaussian-Gaussian KL as the test statistic}
\label{sec:gauss-kl}

For univariate Gaussian distributions with means $\mu_P, \mu_Q$
and variances $\sigma_P^2, \sigma_Q^2$, the forward
Kullback--Leibler divergence has the closed form
\begin{equation}
  \label{eq:gauss-kl}
  \KL\!\bigl(\Nrm(\mu_P, \sigma_P^2)\,\|\,\Nrm(\mu_Q, \sigma_Q^2)\bigr)
  \;=\;
  \tfrac{1}{2}\log\!\frac{\sigma_Q^2}{\sigma_P^2}
  + \frac{\sigma_P^2 + (\mu_P - \mu_Q)^2}{2 \sigma_Q^2}
  - \tfrac{1}{2}.
\end{equation}
Restricting attention to the scalar one-step innovation
(coordinate $0$ of the residual; the only stochastic component
under the rank-one structure of $\hat\Sigma_1$), the directly-fit
and iterated predictive distributions at scale $h$ are univariate
Gaussian:
\begin{align*}
  P_h^{\mathrm{direct}}(x) &= \Nrm\!\bigl((\hat C_h x)_0,\; (\hat\Sigma_h)_{00}\bigr), \\
  P_h^{\mathrm{iter}}(x) &= \Nrm\!\bigl((\hat C_1^{\,h} x)_0,\; (\hat\Sigma_h^{\mathrm{iter}})_{00}\bigr).
\end{align*}
Taking expectation of \eqref{eq:gauss-kl} over the empirical state
distribution restricted to a day-type subgroup $g$, we obtain the
\emph{Gaussian coherence divergence}
\begin{equation}
  \label{eq:gauss-cd}
  \mathrm{kl}(h, g) \;=\;
  \E_{x \sim p_g}\!\left[\,
    \KL\!\bigl(P_h^{\mathrm{direct}}(x) \,\|\, P_h^{\mathrm{iter}}(x)\bigr)
  \,\right].
\end{equation}
Strict-semigroup-consistent factors give $\hat C_h \equiv \hat C_1^{\,h}$
and $\hat\Sigma_h \equiv \hat\Sigma_h^{\mathrm{iter}}$, making
$\mathrm{kl}(h, g) = 0$. Finite-sample noise gives a positive
baseline; the headline test statistic is the ratio against a
synthetic calibration unit:
\begin{equation}
  \label{eq:dratio}
  \mathrm{D\_RATIO}
  \;=\;
  \frac{\mathrm{kl}_{\text{empirical}}(h^{*}, g^{*})}
       {\mathrm{kl}_{\text{synthetic}}(h^{*})},
\end{equation}
where the synthetic system has known Gaussian generator $(L, Q)$
for which (A) holds by construction, $h^{*}$ is a pre-registered
headline horizon, and $g^{*}$ a pre-registered headline subgroup.

\subsection{Gaussian case: what is settled}
\label{sec:settled}

Four campaigns (v1, v2, eigenmodes-v1, P1) ran under the Gaussian
case. Their settled findings are:

\paragraph{(S1) Picture (C) holds for Ontario in the Gaussian case.}
At $(h^{*}, g^{*}) = (12, \mathrm{weekday})$, the 1000-paired-day-bootstrap
gives $\mathrm{D\_RATIO} = 28.96$, $95\%$ CI $[26.82, 31.43]$.
Synthetic baseline $\mathrm{kl}_{\text{synthetic}}(12) \approx 0.0012$
nats (range $0.0011$--$0.0016$ over seven RNG seeds; worst-case
$\mathrm{D\_RATIO} \approx 22$, still above the pre-registered
corroboration bar of $5$). The strict factorization (A) is
empirically inadequate by a margin beyond finite-sample noise.
\emph{Asterisk on the magnitude}: under the corrected (Student-$t$)
kernel family, the analogous statistic $\mathrm{D\_RATIO}_{\Studt}$
at the same headline cell is $2.93$, $95\%$ CI $[2.62, 3.24]$ --- a
$10\times$ shrinkage, below the corroboration bar. The defect exists
(CI clear of $1$, the synthetic floor by construction) but
$\sim 90\%$ of v2's headline magnitude was absorbing the
distributional-class misspecification (S5), not measuring a
genuine cross-scale defect. (S1) is not retracted; its magnitude
is kernel-family-conditional. See \S\ref{sec:q1-settled} (S8) for
the Q1 result that establishes this.

\paragraph{(S2) The defect is diffuse across eigendirections of
$\hat C_1$.}
The dominant-eigenmode share of the defect
$(\hat\Sigma_h - \hat\Sigma_h^{\mathrm{iter}})$ in the eigenbasis of
$\hat C_1$ clusters at $0.50$--$0.58$ across all
$(h, \text{day\_type})$ cells. The single-missing-mode picture
is empirically falsified within the Gaussian setting.

\paragraph{(S3) Genuine day-type asymmetry, not library-noise
inflation.}
The saturday/weekday ratio of $\hat R^{k}(h, \cdot)$ lies in
$[1.26, 2.44]$ across all $(h, k)$ cells; the library-noise
expectation under $1/n$ scaling is $\approx 5$. The asymmetry is
structural.

\paragraph{(S4) Secular drift is the strongest legitimate
calendar signal in PIT residuals.}
On $12{,}000$ post-cutoff PIT values, excluding tautologically-circular
candidates, the year-since-cutoff covariate scores $1.11$ on a
cell-weighted partial-correlation metric; all four seasonal candidates
score in $[0.68, 0.93]$; per-day weekly indicators (except Monday,
borderline) are insignificant.

\paragraph{(S5) The Gaussian forecast distribution itself is wrong
at the marginal level.}
$\E[u] = 0.34$ (uniform: $0.5$); 10-bin $\chi^2 = 22{,}329$ (critical
$16.92$); $67\%$ of PIT values in extreme bins. This is the
empirical observation that prompts the kernel-family generalization
of \S\ref{sec:nongaussian}.

\paragraph{(S6) The Gaussian diffusion ratio runs opposite to seed
prediction.}
$M_2 = (\hat\Sigma_h)_{00} / (\hat\Sigma_h^{\mathrm{iter}})_{00} = 0.73$
at $(h{=}12, \text{weekday})$, against the seed's anticipated
$> 1.5$. At large $h$, the iterated $\hat\Sigma$ over-estimates
because it accumulates beyond the stationary marginal that the
direct $\hat\Sigma_h$ has saturated against.

\paragraph{(S7) The workflow has functioned consistently across
four cycles.}
The arbiter pattern caught one of the three named retraction modes
in each of v1 / eigenmodes-v1 / P1; v2 settled cleanly. The thread
workflow ran end-to-end and surfaced a real workflow gap
(state-advancement provenance) that was patched in the same session.

\paragraph{Summary at the right level of generality.}
Under the Gaussian special case of the framework, picture (C)
holds (S1), the defect has diffuse structure (S2), the asymmetry
is genuine (S3), and secular drift is a measurable component of
the missing content (S4). But (S5) and (S6) suggest the Gaussian
case may not be the right kernel family for the empirical process
in the first place. (S5) is direct: the marginal predictive law
is wrong. (S6) is indirect: the second-moment-only summary
underestimates how iterative variance accumulation diverges from
direct estimation at large $h$.

The settled content of the campaign --- (S1)--(S4) --- is content
\emph{about the Gaussian-projection of the process}, not content
about the process itself in full generality. The framework is
nevertheless validated as a programme of inference: the relaxed
factorization is non-trivial in Ontario data, and its defect has
identifiable structure. Whether the same conclusion holds when
the kernel family is lifted is the next empirical question.

\section{Non-Gaussian extensions}
\label{sec:nongaussian}

The framework of \S\ref{sec:framework} is kernel-family-agnostic.
We here record the structural consequences of lifting the kernel
family and identify what changes vs.\ what stays the same.

\subsection{The decomposition of what depends on Gaussianness}

Of the framework, the Gaussian case used four objects: a kernel
family $\Nrm(\cdot, \cdot)$; a convolution rule \eqref{eq:gauss-compose}
that gives closed-form iteration \eqref{eq:gauss-Ch}--\eqref{eq:gauss-sigma};
a closed-form KL \eqref{eq:gauss-kl}; and a coordinate
standardisation (the z-score transform). We disentangle.

\begin{description}[leftmargin=2em]
\item[(D1) Algebraic framework.] Definitions \ref{def:factorization}
and \ref{def:coherence}, the forward-equation framing
\eqref{eq:forward}, and the strict/relaxed/trivial trichotomy:
\emph{independent} of kernel family.

\item[(D2) Convolution closed form.] Gaussian-family closure
under convolution gives \eqref{eq:gauss-Ch}--\eqref{eq:gauss-sigma}.
Most other families are not closed. For Student-$t$, the
convolution of two Student-$t$ laws is generically not Student-$t$.
For non-parametric families, no closed form exists. In both
non-Gaussian cases, $\hat\kappa_h^{\mathrm{iter}}$ must be computed
\emph{numerically} (by quadrature or by sample-based propagation
of $\hat\kappa_1$). The framework's structural content is unchanged;
only the implementation of \eqref{eq:relax} changes.

\item[(D3) Test statistic.] \eqref{eq:gauss-kl} is closed-form for
Gaussian; for Student-$t$ vs.\ Student-$t$ with matched degrees of
freedom there is a closed-form Student--Student KL; otherwise
\eqref{eq:gauss-cd} is replaced by a sample-based divergence
(Wasserstein-2, energy distance, MMD, or estimated KL via
kernel-density estimation). The choice is one of estimator
efficiency; the structural meaning of the test is preserved.

\item[(D4) Coordinate standardisation.] The z-score transform
$z = (\mathrm{raw} - \mu_{m,h}) / \sigma_{m,h}$ is a deterministic
coordinate change. It removes low-frequency seasonal mean and
scale; it does \emph{not} change the conditional law's shape. The
Gaussian commitment lives not in the z-transform but in the
\emph{predictive distribution} step (the iterated kernel is
Gaussian by assumption). However, the climatology parameters
$(\mu_{m,h}, \sigma_{m,h})$ are themselves estimated under an
implicit Gaussian-second-moment model (the sample mean and
standard deviation per bin). For heavy-tailed innovations, the
standard deviation may not be the right scale parameter, and the
sample mean may differ from the population centre.
\end{description}

\subsection{The general-family test}
\label{sec:general-test}

Generalising \eqref{eq:gauss-cd} to a general kernel family
$\mathfrak{F}$:

\begin{enumerate}[leftmargin=2em]
\item Fit $\hat\kappa_1 \in \mathfrak{F}$ from one-step data
under whichever fitting rule (MLE, OLS, regression with constraints)
respects the family's structure.

\item Compute the iterated kernel $\hat\kappa_h^{\mathrm{iter}}$
either by closed-form convolution (if $\mathfrak{F}$ is closed) or
by numerical propagation (typically: sample $M$ particles from
$\hat\kappa_1(x, \cdot)$; resample; iterate $h$ times; record the
empirical distribution of trajectories).

\item Fit $\hat\kappa_h$ directly from $h$-step data, again under
the family-respecting rule.

\item Compute a divergence
$d_{\mathfrak{F}}\!\bigl(\hat\kappa_h(x, \cdot),\;
\hat\kappa_h^{\mathrm{iter}}(x, \cdot)\bigr)$ for each state $x$
in the empirical distribution. Aggregate by the same subgroup
$g$ stratification.

\item Calibrate against a synthetic system with known generator
in $\mathfrak{F}$ for which strict factorization holds; use the
same divergence procedure to fix the noise floor.

\item Pre-registered test: $\mathrm{D\_RATIO}_{\mathfrak{F}}$ as
in \eqref{eq:dratio}, with thresholds calibrated by the synthetic
noise floor specific to $\mathfrak{F}$.
\end{enumerate}

This is the same six-step procedure as the Gaussian case; only
the family $\mathfrak{F}$ and the divergence $d_{\mathfrak{F}}$
change.

\subsection{The diffusion identity, generalized}
\label{sec:gen-diffusion}

The Gaussian diffusion identity \eqref{eq:lyapunov} has a
generalisation that makes precise which strict-semigroup content
is independent of family. For a general $h$-step kernel with finite
second moments, write the (conditional) drift and the (conditional)
\emph{dispersion} $V_h(x) := \mathrm{Cov}[X' \mid X = x]$ under
$\hat\kappa_h$. Under strict factorization, taking second moments
of the convolution gives the dispersion semigroup identity
\begin{equation}
  \label{eq:disp}
  V_{h_1 + h_2}(x) \;=\; \hat C_{h_2}\, V_{h_1}(x)\, \hat C_{h_2}^{\top}
    \;+\; \E_{y \sim \hat\kappa_{h_1}(x, \cdot)}\!\bigl[\, V_{h_2}(y)\,\bigr].
\end{equation}
For Gaussian kernels with $V_h \equiv \hat\Sigma_h$ constant in
$x$, \eqref{eq:disp} reduces to the iteration rule
$\hat\Sigma_{h_1 + h_2} = \hat C_{h_2} \hat\Sigma_{h_1}
\hat C_{h_2}^{\top} + \hat\Sigma_{h_2}$, which unrolls to
\eqref{eq:gauss-sigma}. For non-Gaussian kernels where $V_h$ may
depend on $x$, the identity does not collapse to closed form, but
it remains a meaningful constraint: the dispersion at scale
$h_1+h_2$ must be the linear-drift-propagated dispersion at scale
$h_1$ plus the average dispersion of the $h_2$-kernel at the
$h_1$-propagated state.

Identity \eqref{eq:disp} is the structural cross-scale
\emph{second-moment} constraint regardless of kernel family.
Higher moments and tail behaviour have analogous identities;
together they exhaust the strict-semigroup content. The
sample-based divergence $d_{\mathfrak{F}}$ implicitly tests all
moments simultaneously; the closed-form Gaussian KL tests the
first two moments specifically.

\subsection{The Student-$t$ extension as a concrete case}
\label{sec:student-t}

The smallest non-Gaussian extension that admits a continuous limit
to Gaussian is the Student-$t$ family. We give the density, the
MLE, the explicit convolution non-closure, and the closed-form
divergence (where one exists) at the same level of detail as
\S\ref{sec:gaussian}.

\paragraph{Density.}
The univariate Student-$t$ location-scale density with degrees of
freedom $\nu \in (0, \infty]$, location $\mu \in \Rb$, and scale
$s > 0$ is
\begin{equation}
  \label{eq:t-density}
  f_{\nu, \mu, s}(y) \;=\;
  \frac{\Gamma\!\bigl(\tfrac{\nu+1}{2}\bigr)}
       {\sqrt{\nu \pi}\, s\, \Gamma\!\bigl(\tfrac{\nu}{2}\bigr)}
  \left(1 + \frac{(y - \mu)^{2}}{\nu\, s^{2}}\right)^{-\frac{\nu+1}{2}}.
\end{equation}
The first two moments are $\E[Y] = \mu$ (for $\nu > 1$) and
$\Var[Y] = s^{2} \cdot \nu / (\nu - 2)$ (for $\nu > 2$;
infinite for $\nu \in (1, 2]$). The Gaussian
$\Nrm(\mu, s^{2})$ is recovered as $\nu \to \infty$. The CDF
$F_{\nu, \mu, s}(y) = \int_{-\infty}^{y} f_{\nu, \mu, s}(y')\, dy'$ has
no elementary form but is available in standard scientific libraries.

The Student-$t$ kernel for the $h$-step conditional law is
\begin{equation}
  \label{eq:t-kernel}
  \hat\kappa_h^{(\Studt)}(x, \cdot) \;=\;
  \Studt_{\nu_h}\!\bigl(\hat C_h x,\; \hat s_h^{2}\bigr),
  \qquad y \sim \hat\kappa_h^{(\Studt)}(x, \cdot)
  \;\Leftrightarrow\;
  y \sim f_{\nu_h,\, \hat C_h x,\, \hat s_h}(\cdot).
\end{equation}
The model has $d^{2}$ drift parameters (one per entry of $\hat C_h$),
one scale parameter $\hat s_h^{2}$, and one degrees-of-freedom
parameter $\nu_h$ per scale $h$.

\paragraph{Maximum likelihood estimation.}
Given $n$ regression pairs $\{(x_i, y_i)\}_{i=1}^{n}$ with $y_i \in \Rb$ the
scalar one-step innovation, the log-likelihood under
$y_i \sim f_{\nu, C x_i, s}$ is
\begin{equation}
  \label{eq:t-loglik}
  \ell(C, s, \nu) \;=\; \sum_{i=1}^{n} \log f_{\nu, C x_i, s}(y_i)
  \;=\; \text{const}(\nu)
    \;-\; \frac{\nu + 1}{2} \sum_{i=1}^{n}
      \log\!\left(1 + \frac{r_i^{2}}{\nu s^{2}}\right)
\end{equation}
with $r_i := y_i - C x_i$ and $\text{const}(\nu)
:= n \log \Gamma(\tfrac{\nu+1}{2}) - n \log \Gamma(\tfrac{\nu}{2})
- n \log(\sqrt{\nu \pi}\, s)$. The MLE alternates between two
sub-problems:

\emph{Regression update} (with $\nu$ fixed). Setting
$\partial \ell / \partial C = 0$ gives the iteratively-reweighted
least-squares (IRLS) normal equations: the weights
$w_i^{(k)} := (\nu + 1) / \bigl(\nu + (r_i^{(k)})^{2}/s^{(k)\, 2}\bigr)$
yield the regression
\begin{equation}
  \label{eq:t-irls}
  C^{(k+1)} \;=\;
  \left(\sum_{i} w_i^{(k)}\, x_i x_i^{\top}\right)^{-1}
  \left(\sum_{i} w_i^{(k)}\, x_i y_i\right),
\end{equation}
iterated with the residuals $r_i^{(k+1)} := y_i - C^{(k+1)} x_i$.
The IRLS converges to a quasi-Newton step of the joint optimisation
in $(C, s)$.

\emph{Scale and degrees-of-freedom updates}. With $C$ held fixed,
$\partial \ell / \partial s^{2} = 0$ gives
\begin{equation}
  \label{eq:t-scale}
  \hat s^{2} \;=\; \frac{1}{n} \sum_{i=1}^{n}
    \frac{(\nu+1)\, r_i^{2}}{\nu + r_i^{2}/s^{2}}
\end{equation}
(implicit; solve by Newton iteration since $s$ appears on both
sides), and the one-dimensional score
$\partial \ell / \partial \nu = 0$ is solved by bracketing on a
log-grid (typical search range $\nu \in [2, 300]$;
boundary failures fall under outcome R-D of Q1).

\emph{Bayesian-regularised alternative}: for finite samples with
small $n$ per day-type, the joint MLE on $(C, s, \nu)$ can be
unstable. A Jeffreys-prior posterior on $\nu$ (with the standard
$\nu^{-1}$ prior, equivalent to a flat prior on $\log \nu$) tames
the boundary behaviour and gives posterior-mean estimates that
are well-defined for any $n$. Pre-registered fallback under R-D.

\paragraph{Convolution non-closure: the explicit statement.}
Let $T_1 \sim \Studt_{\nu}(0, s^{2})$ and $T_2 \sim \Studt_{\nu}(0, s^{2})$
be independent. The density of $T_1 + T_2$ is the convolution
\begin{equation}
  \label{eq:t-convolution}
  f_{T_1 + T_2}(z) \;=\;
  \int_{-\infty}^{\infty}
    f_{\nu, 0, s}(t)\, f_{\nu, 0, s}(z - t)\, dt.
\end{equation}
This integral does not yield a Student-$t$ density. For
$\nu \in \Nb$, it has a closed form expressible via hypergeometric
functions, but the result is itself a non-$t$ distribution. The
tail behaviour is the salient feature: as $|z| \to \infty$, the
density falls polynomially as $|z|^{-(\nu + 1)}$, the same rate as
the original Student-$t$ (each iid summand contributes
multiplicatively to the density, and the slowest-decaying factor
dominates). Thus the \emph{tail index} $\nu$ is preserved under
convolution, but the \emph{shape} of the distribution is not.
In particular, $T_1 + T_2$ is not generally $\Studt_{\nu'}(0, s'^{2})$
for any $(\nu', s')$, and the iterated kernel
$\hat\kappa_h^{(\Studt)\, \mathrm{iter}}$ is outside the parametric
Student-$t$ family for $h \geq 2$.

This algebraic non-closure has the structural consequence we
flagged in \S\ref{sec:drift-vs-rest}: the strict-factorization
constraint (A) is \emph{tighter} in the Student-$t$ case than in
the Gaussian case, because it imposes that $\hat\kappa_{h}^{(\Studt)}$
(directly fit) and $\hat\kappa_h^{(\Studt)\, \mathrm{iter}}$
(iterated) are equal as distributions, even though the latter is
\emph{not Student-$t$}. In the Gaussian case, both sides are
Gaussian by closure and the equality reduces to matching
$(\hat C_h, \hat\Sigma_h)$ to $(\hat C_1^{\, h}, \hat\Sigma_h^{\mathrm{iter}})$.
In the Student-$t$ case, no parameter match suffices; the test must
compare the full distributions.

\paragraph{The directly-fit factor at scale $h$.}
For an independent observation pair $(x_t, y_{t+h})$ where
$y_{t+h}$ is the scalar innovation coordinate at horizon $h$, fit
\eqref{eq:t-irls}--\eqref{eq:t-scale}-and-$\nu$-update on the
sample $\{(x_t, y_{t+h})\}_t$ from the pre-cutoff library. This
yields $(\hat C_h^{(\Studt)}, \hat s_h^{(\Studt)\,2}, \hat\nu_h)$ for
each $h \in \Hcal$. The directly-fit kernel is then
$\hat\kappa_h^{(\Studt)}(x, \cdot) = \Studt_{\hat\nu_h}(\hat C_h^{(\Studt)} x,
\hat s_h^{(\Studt)\,2})$, a closed-form Student-$t$.

\paragraph{The iterated factor at scale $h$.}
Iteration of $\hat\kappa_1^{(\Studt)}$ has no closed form
(\ref{eq:t-convolution}). We compute it by sequential Monte Carlo:
sample $M$ particles $\{u_j^{(1)}\}_{j=1}^{M}$ from $\hat\kappa_1^{(\Studt)}(x, \cdot)$
using the location-scale-degrees-of-freedom representation
$u_j^{(1)} = \hat C_1 x + \hat s_1 T_j$ with $T_j \sim \Studt_{\hat\nu_1}(0, 1)$;
embed each particle by lagged shifts to form an updated state;
iterate $h$ times. The resulting empirical distribution of the
$h$-step coordinate-0 values is the SMC estimate of
$\hat\kappa_h^{(\Studt)\,\mathrm{iter}}(x, \cdot)$, which we
denote $\hat\kappa_h^{(\Studt)\,\mathrm{iter}, M}$ when the particle
count matters. As $M \to \infty$, the SMC estimate converges weakly
to the true iterated distribution.

\paragraph{Closed-form Student-$t$ KL: the matched-$\nu$ case.}
When two Student-$t$ distributions share the same degrees of freedom
$\nu$, the KL divergence has the closed form (after lengthy
calculation; see e.g.\ Bishop, \emph{Pattern Recognition and Machine
Learning}, App.\ B):
\begin{equation}
  \label{eq:t-kl-matched}
  \KL\!\bigl(\Studt_{\nu}(\mu_P, s_P^{2}) \,\|\, \Studt_{\nu}(\mu_Q, s_Q^{2})\bigr)
  \;=\;
  \log\!\frac{s_Q}{s_P}
  + \frac{\nu+1}{2}\, \E_{Y \sim \Studt_{\nu}(\mu_P, s_P^{2})}
    \!\left[\log\!\left(\frac{\nu s_Q^{2} + (Y - \mu_Q)^{2}}
                              {\nu s_P^{2} + (Y - \mu_P)^{2}}\right)\right]
\end{equation}
where the expectation is over a Student-$t$ random variate and
admits an exact representation in terms of digamma and Gauss
hypergeometric ${}_2F_1$ functions, omitted here for brevity. For
the Gaussian limit $\nu \to \infty$, \eqref{eq:t-kl-matched}
recovers \eqref{eq:gauss-kl}.

\paragraph{Mismatched-$\nu$ case: sample-based divergence.}
For Q1, the directly-fit kernel has $\hat\nu_h$ and the iterated
kernel is not Student-$t$ at all (it has tail index $\hat\nu_1$
but non-Student-$t$ shape). The matched-$\nu$ closed form does
not apply. Two practical alternatives:

\emph{(i) Sample-based KL.} Estimate
$\KL(\hat\kappa_h^{(\Studt)} \| \hat\kappa_h^{(\Studt)\,\mathrm{iter},M})$
by
\begin{equation}
  \label{eq:sample-kl}
  \widehat{\KL} \;=\; \frac{1}{N} \sum_{i=1}^{N}
    \log f_h(Y_i)
  \;-\; \frac{1}{N} \sum_{i=1}^{N}
    \log \tilde f_h^{(M)}(Y_i)
\end{equation}
where $Y_i \sim \hat\kappa_h^{(\Studt)}(x, \cdot)$ (drawn from the
fitted Student-$t$), $f_h$ is its closed-form density, and
$\tilde f_h^{(M)}$ is a kernel-density estimate of the iterated
particle ensemble. Bias and variance scale as $O(N^{-1/2})$ and
the KDE bandwidth choice is pre-registered.

\emph{(ii) Wasserstein-2 distance.} Compute
$W_2^{2}\bigl(\hat\kappa_h^{(\Studt)}(x, \cdot),\,
\hat\kappa_h^{(\Studt)\,\mathrm{iter}, M}(x, \cdot)\bigr)
= \inf_{\pi} \int |y_1 - y_2|^{2}\, \pi(dy_1, dy_2)$
between the two univariate distributions. For 1D, $W_2$ has the
explicit formula
\begin{equation}
  \label{eq:w2}
  W_2^{2}(P, Q) \;=\; \int_{0}^{1} \bigl(F_P^{-1}(u) - F_Q^{-1}(u)\bigr)^{2}\, du,
\end{equation}
i.e.\ the $L^2$ distance between quantile functions. The iterated
distribution's quantile function is estimated from the $M$ SMC
particles; the directly-fit distribution's quantile function is
the Student-$t$ inverse CDF. $W_2$ is metric (symmetric and
respects the triangle inequality) and has stable sample-estimator
behaviour even with modest $M, N$.

\paragraph{Coherence divergence under the Student-$t$ kernel.}
The Student-$t$ analogue of the Gaussian coherence divergence
\eqref{eq:gauss-cd} is
\begin{equation}
  \label{eq:t-cd}
  \mathrm{w2}^{(\Studt)}(h, g) \;=\;
  \E_{x \sim p_g}\!\left[
    W_2^{2}\bigl(\hat\kappa_h^{(\Studt)}(x, \cdot),\,
                  \hat\kappa_h^{(\Studt)\,\mathrm{iter}, M}(x, \cdot)\bigr)
  \right]
\end{equation}
with units of (scalar-innovation)$^{2}$, comparable across $h$
and across day-types $g$. The headline test statistic is
\begin{equation}
  \label{eq:dratio-t}
  \mathrm{D\_RATIO}_{\Studt}
  \;=\;
  \frac{\mathrm{w2}^{(\Studt)}_{\text{empirical}}(h^{*}, g^{*})}
       {\mathrm{w2}^{(\Studt)}_{\text{synthetic}}(h^{*})}
\end{equation}
calibrated against a synthetic Student-$t$-VAR system with known
$(L, Q, \nu_{\text{true}})$ for which the strict factorization
holds by construction in the sense that the data are generated by
a single one-step kernel iterated $h$ times.

\paragraph{What the Student-$t$ extension does and does not change.}
The structural framework (Definitions \ref{def:factorization}--%
\ref{def:coherence}, Hypothesis \ref{hyp:working}) is identical
to the Gaussian case --- only \eqref{eq:t-density},
\eqref{eq:t-kernel}, \eqref{eq:t-irls}--\eqref{eq:t-scale},
\eqref{eq:t-convolution}, and the divergence statistic
\eqref{eq:t-cd}--\eqref{eq:dratio-t} replace their Gaussian
counterparts. The dispersion semigroup identity \eqref{eq:disp}
still holds (when $\nu > 2$ so second moments exist) and gives
the same first-and-second-moment cross-scale constraint; what
changes is that the strict factorization additionally imposes
constraints on higher moments and tail behaviour, all of which
the sample-based $W_2$ test \eqref{eq:t-cd} captures at once.

The empirical question --- whether the Gaussian
$\mathrm{D\_RATIO} = 29$ was driven mainly by the kernel-family
mismatch or mainly by the cross-scale defect --- has now been
answered by Q1 of the distributional-class thread (\S\ref{sec:q1-settled}).
The short answer is \emph{both, in a specific ratio}: the family
mismatch accounted for the bulk of the magnitude (the Student-$t$
analogue shrinks roughly $10\times$ to $\mathrm{D\_RATIO}_{\Studt}
= 2.93 \in [2.62, 3.24]$), but a measurable cross-scale defect
remains (CI clear of the synthetic noise floor by construction).
Picture (C) survives the family lift at a much reduced magnitude;
the strict factorization (A) is not recovered. The seed's working
hypothesis is corroborated independent of family, but the lab's
prior reading of v2 as a $29\times$-the-noise-floor effect was
inflated by the wrong kernel family by an order of magnitude.

\subsection{Standardisation under a heavy-tailed kernel}
\label{sec:nongauss-std}

The z-score transform with Gaussian-moment-based
$(\mu_{m,h}, \sigma_{m,h})$ is consistent with a Gaussian kernel
family but not with a Student-$t$ kernel family. For the latter,
the natural standardisation uses Student-$t$ MLE per $(m, h)$ bin:
fit $(\mu_{m,h}, s_{m,h}, \nu_{m,h})$ by maximum likelihood, then
standardise via $z' = (\mathrm{raw} - \mu_{m,h}) / s_{m,h}$ (or
leave the $\nu_{m,h}$ explicit in the kernel). This is a strictly
more general coordinate transform than the Gaussian z-score; in
the limit $\nu_{m,h} \to \infty$ for all bins, the two transforms
agree.

The standardisation question is therefore a function of the kernel
family commitment: choose the family, choose a standardisation
consistent with it. The z-transform is not itself the Gaussian
commitment; rather, the choice of $(\mu_{m,h}, \sigma_{m,h})$ as
the standardisation parameters \emph{is} a partial Gaussian
commitment that is implicit in the current registered estimator.

\section{Q1 of the distributional-class thread: settled}
\label{sec:q1-settled}

The distributional-class thread is rooted in the PIT-calendar-fingerprint
arbiter (S5). Its root node Q1 tested the Student-$t$ extension
against the Gaussian baseline. The experiment ran 2026-05-26; this
section reports the result.

\paragraph{Q1: candidate set (three nested models).}
\begin{itemize}[leftmargin=2em]
\item M0: the registered Gaussian predictor, as null baseline.
\item M1: Student-$t$ kernel with one global $\nu$ per day-type;
location and scale fit by Student-$t$ MLE on the same pre-cutoff
library (\eqref{eq:t-irls}--\eqref{eq:t-scale}); the standardisation
unchanged from M0.
\item M2: full Student-$t$ commitment --- climatology refit under
Student-$t$ MLE per $(m, h)$ bin, plus M1's kernel.
\end{itemize}

\paragraph{Q1: pre-registered outcomes.}
\begin{itemize}[leftmargin=2em]
\item \textbf{R-A:} the marginal PIT under the best of M1, M2 has
$\chi^{2}_{10\text{-bin}} \leq 30$. The Gaussian was the marginal
problem; the heavy-tailed kernel fixes it.
\item \textbf{R-B:} marginal PIT $\chi^{2} \in (30, 500]$.
Student-$t$ helps but does not fully calibrate.
\item \textbf{R-C:} marginal PIT $\chi^{2} > 500$ under both M1
and M2. Student-$t$ does not absorb (S5); the missing content
is not the distributional class.
\item \textbf{R-D:} the MLE fit does not converge or $\nu$
estimates are pathological. Phase~A design fault; redesign.
\end{itemize}

\paragraph{Q1: result.}
With $1000$-paired-day bootstrap CIs:

\begin{center}
\begin{tabular}{lrrr}
\toprule
& point & CI low & CI high \\
\midrule
$\chi^{2}_{M0}$ (Gaussian baseline)        & $22{,}328$ & $19{,}028$ & $25{,}731$ \\
$\chi^{2}_{M1}$ (Student-$t$ kernel only)  & $1{,}371$  & $1{,}069$  & $1{,}817$  \\
$\chi^{2}_{M2}$ (full Student-$t$)         & $1{,}754$  & $1{,}383$  & $2{,}267$  \\
$\mathrm{D\_RATIO}_{\Studt}$ at $(h{=}12, \mathrm{wkdy})$ & $2.93$ & $2.62$ & $3.24$ \\
\bottomrule
\end{tabular}
\end{center}

Empirical $\hat\nu$ per day-type: weekday $4.19$, saturday $4.62$,
sunday $4.61$. The kurtosis $\to \nu$ mapping
$\nu = 4 + 6/\mathrm{excess\_kurt}$ applied to the in-sample
excess kurtosis $24$--$33$ from \texttt{mitacs-rebaseline-facts}
gives $\nu \approx 4.2$--$4.25$; the out-of-sample fits land
within $0.4$ of this.

\paragraph{Q1: verdict.}
The mechanical verdict per the pre-registered criteria is \textbf{R-D}:
$207$ of $288$ M2 climatology bins fitted $\hat\nu$ at the upper
boundary ($300$), tripping the ``$\hat\nu > 200$ pathological'' rule.
Substantively this is the Student-$t$ MLE's correct response when
the bin-level data is approximately Gaussian (CLT-like aggregation
over $\sim 690$ weather samples per $(m, h)$ bin makes the marginal
demand distribution at that month-and-hour approximately Gaussian
even when the conditional residual given recent state is heavy-tailed).
The R-D criterion conflated ``estimator broken'' with ``data
approximately Gaussian at this aggregation level.'' This is the third
occurrence of \emph{pre-registered criterion fires on legitimate-but-
unanticipated data structure} in the workflow's history (after
eigenmodes-v1 and P1); the arbiter pattern handles it the same way
(mechanical verdict honoured, substantive interpretation in
\texttt{verdict\_reasoning}).

The substantive verdict is \textbf{R-C}: $\chi^{2}_{M1} = 1371 > 500$.
Student-$t$ reduces marginal $\chi^{2}$ by $16\times$ but does not
absorb the marginal pathology. The missing content is more than
just the kernel family.

\paragraph{(S8) The Q1 findings, four sub-claims.}
\begin{itemize}[leftmargin=2em]
\item \textbf{(S8a) $\hat\nu$ predictions essentially exact.}
Both forecasters predicted $\nu \approx 4.5$ from the kurtosis
mapping; observed $4.19/4.62/4.61$. The in-sample non-Gaussianity
finding from \texttt{mitacs-rebaseline-facts} holds out-of-sample
on the forecast-grade residual.

\item \textbf{(S8b) Student-$t$ kernel reduces $\chi^{2}$ by $16\times$
but does not fix it.} $22{,}328 \to 1{,}371$. The marginal pathology
has substantial non-distributional content; richer family (mixtures,
non-parametric) or non-distributional causes (mean-trajectory bias,
across-$h$ correlation) remain to be tested.

\item \textbf{(S8c) $\mathrm{D\_RATIO}_{\Studt} = 2.93 \in [2.62,
3.24]$: $10\times$ shrinkage from v2's $28.96$.} Under the corrected
kernel the cross-scale defect at the headline cell sits below the
v2 corroboration bar of $5$. Picture (C) is not refuted (the defect
exists; CI clear of the synthetic noise floor by construction) but
v2's interpretation requires the asterisk on (S1) above:
$\sim 90\%$ of v2's headline magnitude was absorbing distributional-
class misspecification, not measuring a genuine cross-scale defect.
The seed's working hypothesis survives the kernel-family lift at
substantially reduced amplitude.

\item \textbf{(S8d) Heavy-tailedness lives in the $\hat\kappa_1$
innovation, not in the seasonal climatology.} M1 outperforms M2:
adding a Student-$t$ climatology refit \emph{worsens} the marginal
PIT (M2: $\chi^{2} = 1754$ vs M1: $1371$). The bin-level marginal
demand distribution given $(m, h)$ is approximately Gaussian; only
the conditional residual given recent state is heavy-tailed. This
refines where the missing-content programme should look next: not at
the seasonal climatology's distributional form, but at the
conditional residual's structure beyond a single $\nu$ parameter.
\end{itemize}

\paragraph{Forecaster calibration.}
Both forecasters were correct on $\hat\nu$ (the load-bearing
parameter): proponent predicted $4.5$, observed $4.19/4.62/4.61$;
DA agreed with proponent on $\hat\nu$. Both were wrong on the
$\chi^{2}$ magnitude: proponent forecast $\chi^{2}_{M1} = 18$
(observed $1371$; $76\times$ over), DA forecast $250$ ($5.5\times$
over). DA closer but neither close. The substantive verdict R-C
wasn't on either side's slate. The arbiter awarded partial
calibration credit to both: correct mechanism (heavy-tailed kernel
matters), wrong final-outcome label.

\paragraph{Thread state after Q1.}
The distributional-class thread advances from \texttt{planning}
to \texttt{active}. Q1 settled; Q2A (richer family) and Q2B (non-
distributional cause) closed by branching rule because R-D fired
mechanically and the rules send R-D to Q2D (MLE redesign). The
arbiter explicitly flagged this advancement as requiring a thread
amendment before Q2D's phase~A is written: the R-D criterion lumped
the legitimate edge case with broken estimator. A revised R-D
criterion would route Q1 substantively to $\text{R-C} \to \text{Q2A}$.
After amendment, Q2A (richer family --- mixtures, non-parametric)
is the natural next experiment, motivated by (S8b)'s ``Student-$t$
not enough'' finding.

\section{Q2A of the distributional-class thread: settled (R-B2)}
\label{sec:q2a-settled}

Q1's substantive verdict (\S\ref{sec:q1-settled}) was that the
Student-$t$ kernel reduces the marginal $\chi^{2}$ by $16\times$
but leaves $\chi^{2}_{M1} = 1{,}371$ on the table. The natural next
question --- and the proponent's standing read after the Q1 arbiter
flagged R-D as a mechanical artefact --- is whether a richer
distributional family closes the remaining gap. Q2A tests this.
The experiment ran 2026-05-26; this section reports the result.

\paragraph{Q2A: candidate set (three richer-than-Student-$t$ families).}
\begin{itemize}[leftmargin=2em]
\item M3: drift $\hat C_1$ plus a symmetric two-component centred
Gaussian mixture residual law, fitted by EM ($1$ mixing weight $+ 2$
scales per day-type).
\item M4: drift $\hat C_1$ plus a symmetric three-component centred
Gaussian mixture ($2$ mixing weights $+ 3$ scales per day-type).
\item M5: drift $\hat C_1$ plus a non-parametric KDE residual law
(Silverman bandwidth; effectively unbounded shape flexibility on
the $\sim 5{,}000$-sample pre-cutoff library).
\end{itemize}
All three iterate by SMC particle propagation ($M = 200$); only the
one-step residual sampler swaps. The seasonal climatology and drift
are held at the Q1 M1 values.

\paragraph{Q2A: pre-registered cuts.}
Cuts are derived from the Q2A synthetic gate (METHOD/DESIGN-grade
artifact at \path{experiment/results/q2a_synthetic_gate.txt},
\verb|body_sha256| prefix \verb|2bab514e|), a three-variant weight
sweep at $K = 500$ reps per cell on a $4$-point kurtosis sweep
through the production fitters. The variant-max rule gives:
\begin{itemize}[leftmargin=2em]
\item \textbf{R-A2 (corroboration):} $\mathrm{effective\_\chi^{2}}
\leq 228.44$. The richer family closes the gap.
\item \textbf{R-B2 (ambiguous):} $228.44 < \mathrm{effective\_\chi^{2}}
\leq 2{,}284.44$. Richer families help but do not absorb.
\item \textbf{R-C2 (falsification):} $\mathrm{effective\_\chi^{2}}
> 2{,}284.44$. Distributional flexibility above Student-$t$ is
insufficient.
\end{itemize}
$\mathrm{effective\_\chi^{2}} := \min(\chi^{2}_{M3}, \chi^{2}_{M4},
\chi^{2}_{M5})$ on the post-cutoff marginal PIT.

\paragraph{Q2A: result.}
With $1000$-paired-day bootstrap CIs:

\begin{center}
\begin{tabular}{lrrr}
\toprule
& point & CI low & CI high \\
\midrule
$\chi^{2}_{M3}$ (mixture-$2$ Gaussian) & $953.84$  & $695.51$ & $1{,}356.92$ \\
$\chi^{2}_{M4}$ (mixture-$3$ Gaussian) & $965.81$  & $700.89$ & $1{,}386.14$ \\
$\chi^{2}_{M5}$ (KDE residual law)     & $966.60$  & $703.19$ & $1{,}374.12$ \\
$\mathrm{effective\_\chi^{2}}$         & $953.84$  & $685.23$ & $1{,}356.92$ \\
\bottomrule
\end{tabular}
\end{center}

M3 binds; KDE is the worst by $0.8$ $\chi^{2}$ units. The full
M3-to-M5 spread is $13$ $\chi^{2}$ units ($\sim 1.4\%$ of any
single value). $\mathrm{effective\_\chi^{2}} = 953.84
\in [685, 1{,}357]$ sits in the registered R-B2 band: $4.2\times$
above the R-A2 corroboration bar and $2.4\times$ below the R-C2
falsification bar.

\paragraph{Q2A: verdict.}
The mechanical verdict is \textbf{R-B2}, and the substantive verdict
agrees. This is the first arbiter in the workflow's history without
a mechanical/substantive gap: the registered ambiguous band, sized
by the gate's variant-max same-family null, absorbed the realisation
with room to spare on both sides.

\paragraph{(S9) The Q2A findings, three sub-claims.}
\begin{itemize}[leftmargin=2em]
\item \textbf{(S9a) Distributional-class flexibility is empirically
bounded.} The three richer families are $\chi^{2}$-indistinguishable
(spread $13$ units; CIs overlap nearly completely). The simplest
family binds and the most flexible (KDE) is marginally the worst.
Additional shape flexibility above a $2$-component symmetric
mixture buys near-zero marginal calibration on this population.

\item \textbf{(S9b) Day-type asymmetry dominates the pooled
pathology and is family-invariant.} Per-day-type $\chi^{2}$
(recomputed from the result pickle):
\begin{center}
\begin{tabular}{lrrr}
\toprule
& M3 & M4 & M5 \\
\midrule
weekday  & $299.4$  & $271.8$  & $272.2$ \\
saturday & $1{,}184.7$ & $1{,}183.7$ & $1{,}164.3$ \\
sunday   & $565.3$  & $546.7$  & $571.8$ \\
\bottomrule
\end{tabular}
\end{center}
Saturday is $\sim 4\times$ weekday across all three families.
Weekday alone is close to the R-A2 corroboration bar ($228$).
Saturday is $14\%$ of the post-cutoff library but contributes
$\sim 30\%$ of the pooled $\chi^{2}$. The asymmetry is structural
and survives the kernel-family lift, independently corroborating
(S2)/(S3) on weekday/weekend asymmetry being a feature of the
process rather than a library-noise artefact.

\item \textbf{(S9c) The reduction curve asymptotes.} Stacking the
distributional-class lifts:
\begin{center}
\begin{tabular}{lrr}
\toprule
Lift & $\chi^{2}$ & factor \\
\midrule
M0 Gaussian baseline (S5)         & $22{,}329$ & --- \\
M1 Student-$t$ kernel (Q1, S8b)   & $1{,}371$  & $16\times$ \\
Q2A richer-best (M3, S9)          & $954$      & $1.44\times$ \\
\bottomrule
\end{tabular}
\end{center}
Three orders of magnitude M0$\to$M1; one-and-a-half-fold
M1$\to$Q2A. The remaining $\sim 954$ $\chi^{2}$ is unlikely to
yield to further distributional richness. Picture (C) is not
refuted --- the defect persists with CI clear of $1$ per (S8c) ---
but (S9) tightens (S8c)'s qualifier: the kernel-family-lift
programme has been substantially exhausted on this empirical
process. Further reduction requires non-distributional content
(state-dependent mixtures, hybrid distributional/calendar models,
or non-distributional drivers such as iterated-mean bias, seam
artefacts, or secular drift).
\end{itemize}

\paragraph{Forecaster calibration.}
Proponent forecast R-B2 with $\mathrm{effective\_\chi^{2}}$
central estimate $550$ and CI $[280, 900]$; KDE-preferred. Devil's
advocate forecast R-A2 with central $180$ and CI $[80, 360]$;
KDE-preferred. Realised $953.84$, M3-preferred. Proponent got the
outcome category exactly right and the qualitative shape (richer
families help but do not close the gap) exactly right; central
forecast overshot the realisation by $\sim 5\%$ above the CI upper
bound and the KDE-preferred sub-prediction was wrong (all three
families tied). DA was wrong on outcome, magnitude ($5\times$),
CI, and family ranking; DA's own ``what would change my mind''
clause named $\mathrm{effective\_\chi^{2}} > 1{,}000$ as the
condition for conceding, and the realisation sits just below
that threshold but $2.65\times$ above DA's CI ceiling, which counts
as concession on DA's own terms. Proponent decisively closer.

The proponent's contingency clause ``if all three families produce
nearly identical $\chi^{2}$ within bootstrap CI overlap, the
distributional class beyond Student-$t$ is irrelevant and the
residual gap is purely non-distributional'' (item~$4$ of their
\texttt{what\_would\_change\_my\_mind}) materialised verbatim. The
substantive reading of R-B2 is therefore stronger than the
mechanical statement: not just ``richer families help but do not
fully calibrate,'' but ``additional distributional flexibility
above a $2$-component symmetric mixture buys negligibly little.''

\paragraph{Thread state after Q2A.}
Q2A.status $\to$ \texttt{settled}; the thread's
\texttt{branching\_rules} table has no Q2A entry, so mechanical
state advancement is unavailable and the thread transitions to
\texttt{pending\_amendment}. The arbiter flagged three amendment
candidates for the analyst (not the thread-coordinator): reopen Q2B
(non-distributional content); define a new Q2C with state-dependent
or hybrid distributional/calendar models; or declare the
distributional-class thread exhausted and spawn a non-distributional
thread on the remaining $\sim 954$ $\chi^{2}$. (S9c)'s asymptote
reading favours the third.

\section{Q2B of the distributional-class thread: settled (R-C3)}
\label{sec:q2b-settled}

Q2A's substantive verdict (\S\ref{sec:q2a-settled}) bounded the
distributional-class lift: three richer-than-Student-$t$ families
are $\chi^{2}$-indistinguishable, the simplest binds, and the
remaining $\sim 954$ $\chi^{2}$ does not yield to further
distributional richness. Q2B is the registered ablation experiment
that decomposes that remaining $\chi^{2}$ into three pre-registered
non-distributional mechanisms. It ran 2026-05-26; this section
reports the result.

\paragraph{Q2B: candidate set (three non-distributional mechanisms).}
Each mechanism is a single-target ablation fix on the iterated
$\hat\kappa_{1}$ predictor at the M3 (mixture-$2$ Gaussian)
baseline; only that piece changes:
\begin{itemize}[leftmargin=2em]
\item Mech~1 (\texttt{mean\_bias}): add the empirical per-$(h,
\mathrm{day\_type})$ mean bias to each SMC particle's $z$ value
before marginal PIT. Targets clim-gap-nonstationary's MW-space
$+1000$~MW drift after $z$-space translation.
\item Mech~2 (\texttt{variance\_div}): one-sided variance cap.
Rescale particles toward the ensemble mean if $\sigma^{2}_{\mathrm{iter}}
> \sigma^{2}_{\mathrm{empirical}}$. Targets putative iterated
over-dispersion.
\item Mech~3 (\texttt{seam}): restrict PIT aggregation domain to
exclude seam-domain rows. \textbf{Registered null} (production
iteration already excludes seam-crossing rows via
\verb|_complete_delivery_days|).
\end{itemize}

\paragraph{Q2B: pre-registered cuts.}
The verdict metric is the explained-share of the gap between the
M3 baseline and Q2A's R-A2 corroboration cut:
\[
\mathrm{explained\_share}[m] :=
\frac{\chi^{2}_{\mathrm{baseline\_M3}} - \chi^{2}_{\mathrm{after\_fix}}[m]}
{\chi^{2}_{\mathrm{baseline\_M3}} - \mathrm{R\text{-}A2\_cut}}
\]
with $\chi^{2}_{\mathrm{baseline\_M3}} = 953.84$ (Q2A SETTLED),
$\mathrm{R\text{-}A2\_cut} = 228.44$, and frozen denominator
$725.40$. The shape outcomes are:
\begin{itemize}[leftmargin=2em]
\item \textbf{R-A3 (dominance):} $\max_{m}
\mathrm{explained\_share}[m] > 0.50$ AND $> 2\times$ second
share. A single mechanism captures the bulk.
\item \textbf{R-B3 (mixed):} not R-A3 and not R-C3 (at least one
share $\geq 0.15$). Mechanisms contribute jointly.
\item \textbf{R-C3 (inadequate):} $\max_{m}
\mathrm{explained\_share}[m] < 0.15$. No pre-registered
mechanism captures the gap.
\end{itemize}

\paragraph{Q2B: result.}
With $1000$-paired-day bootstrap CIs:

\begin{center}
\begin{tabular}{lrrr}
\toprule
& point & CI low & CI high \\
\midrule
$\chi^{2}_{\mathrm{baseline\_M3}}$         & $961.96$  & $698.61$  & $1{,}375.68$ \\
$\chi^{2}_{\mathrm{after\_mean\_bias}}$    & $875.04$  & $627.65$  & $1{,}213.67$ \\
$\chi^{2}_{\mathrm{after\_variance\_cap}}$ & $1{,}036.15$ & $750.79$  & $1{,}512.26$ \\
$\chi^{2}_{\mathrm{after\_seam\_excl}}$    & $1{,}054.59$ & $790.12$  & $1{,}464.07$ \\
\midrule
$\mathrm{share}[\mathtt{mean\_bias}]$       & $+0.109$ & $-0.393$ & $+0.735$ \\
$\mathrm{share}[\mathtt{variance\_div}]$    & $-0.114$ & $-0.217$ & $-0.013$ \\
$\mathrm{share}[\mathtt{seam}]$             & $-0.139$ & $-0.169$ & $-0.088$ \\
\bottomrule
\end{tabular}
\end{center}

Self-consistency baseline reproduces Q2A's settled M3 to within
$0.85\%$ ($|\Delta| = 8.12$ vs.\ tolerance $50$; Q2A CI overlap
clear). Pipeline integrity confirmed.

\paragraph{Q2B: verdict.}
The mechanical verdict is \textbf{R-C3} (max share $0.109$ on
\texttt{mean\_bias}, well below the $0.15$ null floor), and the
substantive verdict agrees. Q2B is the second arbiter in the
workflow's history without a mechanical/substantive gap; but here
the substantive picture is \emph{stronger} than the mechanical
reading, not weaker, because two of three shares are
sign-determined NEGATIVE at $95\%$ --- the corresponding ablation
fixes actively made $\chi^{2}$ worse.

\paragraph{(S10) The Q2B findings, three sub-claims.}
\begin{itemize}[leftmargin=2em]
\item \textbf{(S10a) All three pre-registered mechanism candidates
are empirically inadequate.} Point shares
$(+0.109, -0.114, -0.139)$ all below the $0.15$ null floor.
The pre-registered mechanism set (mean-bias absorption, one-sided
variance cap, seam-domain exclusion) does \emph{not} capture the
non-distributional structure of the $\sim 954$ $\chi^{2}$ residual.

\item \textbf{(S10b) Two mechanisms are mis-specified at the
construction level, not merely inadequate.} \texttt{variance\_div}
and \texttt{seam} have negative shares sign-determined at $95\%$
(CIs $[-0.22, -0.01]$ and $[-0.17, -0.09]$ respectively). The
ablation fixes did not just fail to improve calibration; they
degraded it. The mechanistic reason for \texttt{variance\_div}:
$\sigma^{2}_{\mathrm{iter}} < \sigma^{2}_{\mathrm{empirical}}$
at virtually every $(h, \mathrm{day\_type})$ cell --- the SMC
ensemble UNDER-disperses relative to actuals everywhere. The
one-sided cap is wrong-directional: designed for over-dispersion,
applied in an under-dispersion regime. For \texttt{seam}: the
registered null fires with sign because layering an additional
row-exclusion on production output (which already handles the
seam correctly) removes valid PIT samples and adds aggregation
noise.

\item \textbf{(S10c) The under-dispersion signature is the
substantive positive content.} The SMC ensemble variance is less
than the empirical variance at nearly every cell across all three
day-types (only a minority of saturday cells show
$\sigma^{2}_{\mathrm{iter}} \geq \sigma^{2}_{\mathrm{empirical}}$).
Under-dispersion of an iterated forecast is the canonical signature
of missing exogenous content: the predictor's predictive
distribution is too narrow because the conditioning sub-$\sigma$-algebra
is too coarse. The mech~$1$ contribution ($+0.11$, the only
positive share) is directionally correct but bounded; richer
subtraction designs (state-dependent, weather-conditioned,
secular-drift-conditioned) may extract more, but the aggregate
per-$(h, \mathrm{day\_type})$ ceiling at $\sim 0.11$ is a
substantive upper bound on what aggregate mean-bias ablation
alone can deliver.
\end{itemize}

\paragraph{Forecaster calibration.}
Proponent forecast R-A3 with \texttt{mean\_bias} dominance at
share $0.75$ (CI $[0.50, 0.95]$), MW-space-to-$z$-space anchored.
Devil's advocate forecast R-B3 with \texttt{mean\_bias} share
$0.40$ (CI $[0.25, 0.60]$), $z$-space attenuation anchored. Both
wrong on outcome category; DA materially closer on the
\texttt{mean\_bias} magnitude (residual $0.29$ vs.\ proponent
$0.64$). Both shared a structural blind spot: neither anticipated
that any mechanism ablation could actively degrade calibration.
The proponent's explicit \texttt{what\_would\_change\_my\_mind}
condition $\#2$ (``If R-C3 fires, none of the three pre-registered
mechanisms is the right level of decomposition'') materialised
verbatim. This is the fourth named-failure-mode instance in the
workflow (after Q1 R-D-vs-R-C, eigenmodes-v1 O$2$-vs-O$3$, and
P1 mechanical-AMBIGUOUS-vs-tautology); Q2B differs in direction
--- the substantive verdict is \emph{stronger} than the mechanical
reading, not different from it.

\paragraph{Framework-level interpretation: the identifiability
obstruction.}
The conjunction of (S9) and (S10) admits a sharper reading under
the disintegration-as-diagnostic frame parked
2026-05-26 in the theory-side programme. Decompose the risk of any
estimator built from a conditional kernel $\kappa_{Q}$ as
\[
\mathrm{Risk}(\mathrm{estimator}) \;=\; R^{*}(Q) \;+\; \text{model-class gap} \;+\; \text{estimation error},
\]
where $R^{*}(Q)$ is the conditional Bayes risk at the
conditioning $\sigma$-algebra $Q$. The framework offers two
levers for reducing risk: (1) \emph{kernel refinement} reduces
the model-class gap at fixed $Q$; (2) \emph{$\sigma$-algebra
enrichment} reduces $R^{*}(Q)$ itself by conditioning on more
information ($Q \to Q' \supset Q$). (S9) was a properly-posed
Lever~1 test: three richer-than-Student-$t$ families at the same
$Q$, bounded at $1.44\times$ marginal return. (S10) was \emph{not}
a properly-posed Lever~2 test: the three mechanisms (mean-bias
subtraction, variance cap, seam exclusion) are post-hoc fixes on
the same $Q$, not enrichments. The empirical conjunction --- Lever~1
plateaus and naive same-$Q$ manipulations fail with two
mis-specified --- is consistent with \emph{either} regime~(a)
intrinsic structured noise (the conditional law $\kappa_{Q}(\cdot \mid q)$
is heavy-tailed for intrinsic reasons; $R^{*}(Q)$ is the floor)
\emph{or} regime~(b) noisy structure with latent $Z \notin \sigma(Q)$
orthogonal to the current conditioning (the observed pathology
is a mixture artefact and $R^{*}(Q') < R^{*}(Q)$ for $Q' = Q \vee \sigma(Z)$).

That these two regimes produce identical $\kappa_{Q}$ and are
\emph{provably} indistinguishable from $\kappa_{Q}$ alone is
published theorem in multiple traditions: Bergna et al.
(\texttt{arXiv:2605.06413}, $2026$, Prop.~$1$) for the aleatoric--epistemic
ambiguity, Heckman--Singer (\emph{Econometrica}, $1984$) for
duration models, Allahverdyan (\texttt{arXiv:2002.07884}, $2020$)
for mixture models, and Andrews--Mallows (\emph{JRSSB}, $1974$)
for scale-mixture transparency. The memo's contribution is not
the obstruction theorem; it is applied evidence that (S9) and
(S10) jointly map the boundary of what the distributional-class
axis can deliver on this empirical process, and that the natural
next move follows the framework's prescribed resolution paths:
(i) candidate-$Z$ testing (temperature, secular drift via
\texttt{year\_since\_cutoff} per (S5), refined day-type partition);
(ii) stratified covariate-dependent two-component mixture
decomposition via NPMLEmix (Deb, Saha, Guntuboyina, Sen, JASA
2022; arXiv:1810.07897) as a $\sigma$-algebra adequacy
diagnostic;\footnote{The disintegration note's Phase-2 audit row
listed ``Patra--Sen conditional diagnostic'' as a novel applied
contribution; a second-pass literature scout (2026-05-26;
\path{notes/literature/2026-05-26_stratified-patra-sen-scout.md})
falsified this, finding that Deb et al.\ (2022, by Sen of
Patra--Sen 2016) provide the strictly more general joint NPMLE
of $\pi^*(x)$ with a formal distance-covariance adequacy test in
their Section~6. The methodology is canonical; the diagnostic
value is unchanged.} (iii) synthetic Gaussian-DGP calibration to
anchor the residual-$\chi^{2}$ floor against a known reference.

\paragraph{Thread state after Q2B.}
Q2B.status $\to$ \texttt{settled}; the thread's
\texttt{branching\_rules} table has no Q2B entry, so mechanical
state advancement is unavailable and the thread transitions to
\texttt{pending\_amendment}. The substantive recommendation,
under the framework reframe above, is to declare the
distributional-class thread \texttt{exhausted} on its
pre-registered scope (distributional-class axis only) and spawn
a fresh thread on the three resolution paths. The arbiter notes
this is one of two legitimate moves under the workflow (the other
being a mechanism-redesign amendment); the framework reframe
favours the spawn-new-thread move because the three resolution
paths are not mechanism-redesign within the same lever.

\section{P1 of the resolution-paths thread: settled (R-A landslide)}
\label{sec:p1-settled}

(S9) and (S10) jointly bound the distributional-class lever and map
its boundary to the disintegration-as-diagnostic identifiability
obstruction. The natural next move follows the framework's three
prescribed resolution paths (\S\ref{sec:q2b-settled}, last paragraph):
(i) candidate-$Z$ testing, (ii) stratified covariate-dependent
two-component mixture decomposition via NPMLEmix per Deb et al.\
(2022), (iii) synthetic Gaussian-DGP calibration. A new pre-registered
thread (\texttt{resolution-paths-thread}) was opened 2026-05-27 to
test these. Its root node P1 is the diagnostic that does \emph{not}
require auxiliary data: a per-stratum application of the
Patra--Sen (2016) two-component non-Gaussian mixture estimator
$\hat\alpha_L$ to Q1's M1 PIT residuals, stratified along five
candidate axes. If $\hat\alpha_L$ varies materially across strata
on any axis, that axis names where in the existing time-indexed
$\sigma$-algebra the missing-$Z$ structure is concentrated. The
experiment ran 2026-05-27; this section reports the result.

\paragraph{P1: candidate set (five conditioning axes).}
\begin{itemize}[leftmargin=2em]
\item \texttt{hour\_of\_week}: 8 bins of 21 hours, indexed
\(\lfloor (\mathrm{weekday}\cdot 24 + (h{-}1))/21 \rfloor\),
\(\mathrm{weekday}\in\{0,\ldots,6\}\), \(h\in\{1,\ldots,24\}\).
\item \texttt{time\_of\_day}: 4 bins (overnight, morning ramp,
afternoon, evening), per
\path{experiment/p1_patra_sen/__main__.py:81-96}.
\item \texttt{season}: 3 bins (winter Nov--Mar, summer Jun--Aug,
shoulder Apr--May+Sep--Oct).
\item \texttt{day\_type}: 3 bins (weekday, saturday, sunday).
\item \texttt{demand\_quantile}: terciles of the post-cutoff
demand distribution.
\end{itemize}

\paragraph{P1: pre-registered outcomes.}
The verdict statistic is
\(\mathrm{range}_{\mathrm{axis}} = \max_{s\in\mathrm{axis}}\hat\alpha_L(s) -
\min_{s\in\mathrm{axis}}\hat\alpha_L(s)\), maximised over the five
axes, with $1000$-paired-day bootstrap CIs:
\begin{itemize}[leftmargin=2em]
\item \textbf{R-A:} $\max_{\mathrm{axis}}\mathrm{range} > 0.20$. Missing
$Z$ detectable from current data; the binding axis names where
$Z$ is concentrated.
\item \textbf{R-B:} $\max_{\mathrm{axis}}\mathrm{range} \in (0.10, 0.20]$.
Marginal variation; $Z$ hint present but not localised.
\item \textbf{R-C:} $\max_{\mathrm{axis}}\mathrm{range} < 0.10$.
Constant across all strata; $Z$ orthogonal to current
time-indexed conditioning.
\end{itemize}

\paragraph{P1: methodology amendment (pre-data, carve-out).}
The original phase~A specified $F_b = N(0, 1)$ on standardised
residuals $u_{\mathrm{std}} = (z_{\mathrm{actual}} -
z_{\mathrm{iter}})/s_{\mathrm{iter}}$. This was wrong: M1 is a
Student-$t$ kernel, so under correct calibration $u_{\mathrm{std}}$
is Student-$t(\nu)$, not standard normal. Patra--Sen with
$F_b = N(0, 1)$ on a correctly-calibrated Student-$t$ sample returns
a large positive $\hat\alpha_L$ by construction --- not a missing-$Z$
signal. Caught at the orient step before any Ontario data was
touched; re-stamped under the preregister agent's
\emph{pre-run methodology amendment} carve-out (commit
\texttt{99a8122}). Corrected: $F_b = \mathrm{Uniform}[0,1]$ on
$u_{\mathrm{PIT}}$ (Q1's already-emitted PIT values). Patra--Sen
2016 Theorem~1 (monotone invariance of $\alpha_0$) gives the
equivalence; this matches Q1's measurement convention. The
range-shape forecasts are preserved under monotone reparameterisation,
so proponent/DA forecasts stayed forecast-consistent.

\paragraph{P1: result.}
With $1000$-paired-day bootstrap CIs, per-axis ranges (point [CI]),
ordered:

\begin{center}
\begin{tabular}{lrrr}
\toprule
axis & point & CI low & CI high \\
\midrule
\texttt{hour\_of\_week}    & $0.5750$ & $0.4475$ & $0.6475$ \\
\texttt{time\_of\_day}     & $0.4850$ & $0.3850$ & $0.5775$ \\
\texttt{day\_type}         & $0.3450$ & $0.2500$ & $0.4300$ \\
\texttt{season}            & $0.1350$ & $0.0550$ & $0.2400$ \\
\texttt{demand\_quantile}  & $0.0750$ & $0.0225$ & $0.1725$ \\
\bottomrule
\end{tabular}
\end{center}

Only \texttt{demand\_quantile} has a CI overlapping the R-C cut at
$0.10$; the other four carry localising signal. The binding axis
(\texttt{hour\_of\_week}) splits into a clear hot/cold pattern:
hot zone (bins 5/6/7, Fri 10:00 -- Sun 23:00) at
$\hat\alpha_L \in \{0.55, 0.61, 0.55\}$; cold zone (bin 0, Mon
01:00 -- Mon 21:00) at $\hat\alpha_L = 0.035$ (essentially Gaussian).

\paragraph{P1: verdict.}
The mechanical verdict is \textbf{R-A at landslide strength}
($\max_{\mathrm{axis}}\mathrm{range} = 0.575$, lower CI bound
$0.4475$ clears the R-A cut $0.20$ by more than $2\times$). The
substantive verdict agrees on category and adds: the missing-$Z$
structure is \emph{cyclical/social-weekly, not seasonal/
climatological}. Season ranks fourth of five axes at $0.135$,
barely above the R-A cut; the prior-supported temperature candidate
(\path{mitacs-clim-gap-nonstationary}) takes a substantive downward
update on its localising-axis claim (temperature remains candidate
for the marginal pathology, but in post-cutoff PIT residuals the
localising signal sits on a weekly/social grid, not climatology).

\paragraph{(S11) The P1 findings, three sub-claims.}
\begin{itemize}[leftmargin=2em]
\item \textbf{(S11a) The missing-$Z$ structure is cyclical/social-
weekly, NOT seasonal/climatological.} Hot zone is the weekend window
(Fri evening through Sunday); cold zone is Monday morning.
\texttt{hour\_of\_week} range $0.575$ binds; \texttt{season} range
$0.135$ is fourth of five. Temperature prior takes a substantive
hit on the localising-axis claim.

\item \textbf{(S11b) Day-type ranking corroborates two distinct
prior estimators.} Saturday $\hat\alpha_L = 0.6475 >$ sunday
$0.4850 >$ weekday $0.3025$. The Saturday-$4\times$-weekday
asymmetry from Q2A Finding~$2$ (S9) and the weekday/weekend
structural asymmetry from eigenmodes-v1 (S2, S3) are corroborated
from a distinct estimator family ($\hat\alpha_L$ rather than
$\chi^2$).

\item \textbf{(S11c) Patra--Sen \emph{per-stratum} application is
a useful applied diagnostic for $\sigma$-algebra adequacy.} On
post-cutoff data, four of five axes carry detectable localising
signal at $95\%$. The diagnostic does not require auxiliary data
and does not assume a parametric mixture; it tests whether the
$\sigma$-algebra of the predictor is fine enough to absorb the
local non-Gaussian fraction the data carry.
\end{itemize}

\paragraph{Forecaster calibration.}
Proponent forecast R-A with binding axis \texttt{season}, range
$0.30$ (CI $[0.18, 0.45]$). Mechanical category correct
(\textbf{PROPONENT-CONFIRMED} on outcome); within-bucket
substantive forecasts $95\%$-CI-falsified on three of four
sub-quantities (binding axis is \texttt{hour\_of\_week} not
\texttt{season}; range $0.575$ overshoots upper CI $0.45$ by
$0.125$; $\hat\alpha_L$ marginal $0.4125$ undershoots lower CI).
Devil's advocate R-C decisively falsified ($4$ of $5$ per-axis
CIs above the $0.10$ R-C cut); DA's own ``what would change my
mind'' condition $\#2$ (R-A for \texttt{hour\_of\_week} with range
$> 0.25$) materialised verbatim at $0.575$. P1 is the fourth
named-failure-mode instance (criterion-vs-legitimate-data-structure)
in a new orthogonal shape: mechanical matches proponent's category,
within-bucket detail matches neither side.

\paragraph{Thread state after P1.}
P1.status $\to$ \texttt{settled} R-A; thread branching rule
routes R-A to Q1A (re-fit M3 with the binding-axis $Z$ added) and
closes Q1C (the temperature-fallback sibling, since R-A fired).
Thread state \texttt{active}, current\_node = Q1A.

\section{Q1A of the resolution-paths thread: settled (R-B1A, NEAR-TIE)}
\label{sec:q1a-settled}

P1 named \texttt{hour\_of\_week} as the binding localising axis;
Q1A is the registered $\chi^2$-side test of whether enriching M3's
conditioning $\sigma$-algebra with that $Z$ closes the marginal-PIT
gap. The experiment ran 2026-05-27; this section reports the result.

\paragraph{Q1A: candidate set (head-to-head $Z$ constructions).}
\begin{itemize}[leftmargin=2em]
\item $Z_b$: binary weekend indicator ($\mathbb{1}[\mathrm{hour\_of\_
week\_bin} \ge 5]$). Parsimony baseline; the simplest construction
that captures the hot-zone/cold-zone split P1 found.
\item $Z_c$: $8$-level categorical (the full \texttt{hour\_of\_week}
bin index $\in \{0, \ldots, 7\}$). Maximum $\sigma$-algebra
refinement on the axis.
\end{itemize}
Both refit M3 (mixture-$2$ Gaussian) per $(\mathrm{day\_type}, Z)$
cell on the same pre-cutoff library Q2A used; SMC particle
propagation per cell.

\paragraph{Q1A: pre-registered outcomes.}
The verdict statistic is
$\mathrm{better}\_\chi^2 = \min(\chi^2_{Z_b}, \chi^2_{Z_c})$ on
post-cutoff marginal PIT, with the Q2A-derived cuts:
\begin{itemize}[leftmargin=2em]
\item \textbf{R-A1A:} $\mathrm{better}\_\chi^2 \le 228$ (Q2A's
R-A2 corroboration cut). $Z$ closes the gap; the specific $Z$ is
identified.
\item \textbf{R-B1A:} $\mathrm{better}\_\chi^2 \in (228, 2284]$
(Q2A's R-B2 band). $Z$ helps but does not close; multi-$Z$ needed.
\item \textbf{R-C1A:} $\mathrm{better}\_\chi^2 > 2284$. $Z$
doesn't help; P1's signal was a false positive or the
localisation missed the actual $Z$.
\end{itemize}

\paragraph{Q1A: result.}
With $1000$-paired-day bootstrap CIs:

\begin{center}
\begin{tabular}{lrrr}
\toprule
& point & CI low & CI high \\
\midrule
$\chi^2_{Z_b}$ (binary)              & $986.02$ & $678.23$ & $1{,}381.66$ \\
$\chi^2_{Z_c}$ ($8$-level)           & $952.52$ & $626.97$ & $1{,}373.91$ \\
$\mathrm{better}\_\chi^2$            & $952.52$ & $625.34$ & $1{,}342.33$ \\
$\chi^2_{\mathrm{diff}} = Z_b - Z_c$ & $+33.50$ & $-93.16$ & $+142.64$ \\
\bottomrule
\end{tabular}
\end{center}

Paired SE on the difference: $56.84$; $|\chi^2_{\mathrm{diff}}|/
\mathrm{SE} = 0.59 < 1$, so the two $Z$ constructions are
statistically indistinguishable on the verdict-metric (TRUE NEAR-TIE;
$Z_c$ wins at the point estimate but not at $95\%$). The
reproduction check (M3 no-$Z$ vs Q2A's settled M3) is EXACT:
$\chi^2_{\mathrm{no\text{-}Z}} = 953.84$, Q2A settled $= 953.84$,
$|\mathrm{diff}| = 0.0$. The seed-equality clause holds end-to-end;
$Z$-conditioning $\chi^2$ values are directly comparable to Q2A's
baseline.

\paragraph{Q1A: verdict.}
Mechanical: \textbf{R-B1A} (better $952.52$ in $(228, 2284]$),
with NEAR-TIE sub-finding ($Z_c$ wins at point, indistinguishable
at $95\%$). Substantively the result is stronger than the
mechanical category suggests: $Z$ conditioning shaved $1.33$
$\chi^2$ units off Q2A's M3 = $953.84$ baseline, which is
$\mathbf{0.18\%}$ of the $725.84$ gap-to-close to R-A1A = $228$.
The \texttt{hour\_of\_week} $\sigma$-algebra enrichment delivers
essentially no marginal reduction beyond bootstrap noise.

\paragraph{(S12) The Q1A findings, three sub-claims.}
\begin{itemize}[leftmargin=2em]
\item \textbf{(S12a) Per-cell decomposition exactly corroborates
P1.} The saturday $\mathrm{bin}_6$ cell ($Z_c = 6$, the
\texttt{hour\_of\_week} hot zone of Sat 17:00 -- Sun 13:00)
contributes $\chi^2 = 1634$, single cell exceeding the pooled
marginal $952.5$ by $70\%$. This is exactly the P1 stratum at
$\hat\alpha_L = 0.61$. \textbf{The contamination is real and
localised where P1 said it would be.}

\item \textbf{(S12b) The mixture-$2$ Gaussian family is at
capacity within strata.} Per-cell stratification cannot reduce
saturday $\mathrm{bin}_6$'s $\chi^2 = 1634$ below $228$ because
M3's mixture-$2$ Gaussian is not flexible enough to absorb the
residual law there. Combined with (S9)'s finding that M3/M4/M5
distributional flexibility is $\chi^2$-indistinguishable, no
amount of within-$Z$ same-class substitution closes this cell.
The hot-zone non-Gaussianity P1 detected is a family-capacity
bound on Q1A's mechanism, not a $Z$-identification failure.

\item \textbf{(S12c) Three within-data levers are now empirically
bounded.} Joint (S9) + (S10) + (S11) + (S12) maps:
distributional-class (Q2A: $1.44\times$ marginal, asymptotes),
same-$Q$ mechanism ablation (Q2B: all three below null floor;
two sign-determined negative), \texttt{hour\_of\_week} stratification
(Q1A: $0.18\%$ of gap). All three are bounded. Resolution requires
\emph{auxiliary information beyond the time index}; the
identifiability obstruction (Bergna et al.\ 2026 Prop.~1;
Heckman--Singer 1984; Allahverdyan 2020) is now corroborated from
the $\chi^2$ side in addition to the Patra--Sen $\hat\alpha_L$
side.
\end{itemize}

\paragraph{Forecaster calibration.}
Devil's advocate forecast R-B1A with $Z_c$ point chi$^2 = 950$
(CI $[800, 1100]$, modelled on Q2A's M3 carryforward at $953.84$
minus a few $\chi^2$). Realised $Z_c = 952.52$, residual
$\mathbf{0.3\%}$ on the within-bucket point forecast --- the
strongest single calibration data point in the workflow's history.
Proponent forecast R-A1A with $\mathrm{min}\_\chi^2 = 200$ (CI
$[130, 320]$), decisively falsified ($4.8\times$ overshoot of
realised lower CI $625$). Within-bucket detail: DA predicted $Z_b$
wins by $-150$ (realised $Z_c$ wins by $+33.5$, sign reversed but
magnitude $\sim 1\sigma$-of-diff scale correct); proponent
predicted $Z_c$ wins by $+220$ (realised $+33.5$, one-sixth of
proponent's magnitude). Fifth named-failure-mode instance in a
new shape: \textbf{DA-CONFIRMED} on outcome category with
essentially exact within-bucket point forecast but wrong
within-bucket winner direction; proponent fully wrong on category.
The mechanical verdict and substantive verdict align on R-B1A;
the substantive content sharpens the framework reading (within-data
levers exhausted) rather than pulling the verdict toward a
neighbouring category.

\paragraph{Thread state after Q1A.}
Q1A.status $\to$ \texttt{settled} R-B1A; thread branching rule
routes R-B1A to Q1B (re-fit M3 with a second $Z$ added on top of
$Z_c$, drawn from P1's second-strongest axis \texttt{time\_of\_day}).
Carryforward to Q1B: $Z_c$ is the parsimony-preferred point-estimate
winner with the explicit NEAR-TIE caveat that $Z_b$ vs $Z_c$ is
not statistically discernible. Thread state \texttt{active},
current\_node = Q1B. An advisory pre-flag is logged on the thread's
state history: given Q1A's $0.18\%$-of-gap reading on the
\emph{binding} axis, Q1B (\texttt{time\_of\_day}, P1 range $0.485$
vs \texttt{hour\_of\_week} $0.575$) faces the same skepticism with
weaker prior support; if Q1B also fires R-B1B/R-C1B at similar
magnitudes, time-indexed $\sigma$-algebra resolution paths are
empirically exhausted and the thread requires an amendment for
temperature/exogenous-$Z$. The advisory is background; Q1B's
verdict is not pre-judged.

\section{Q1B of the resolution-paths thread: settled (R-B1B, substantive FAVORABLE)}
\label{sec:q1b-settled}

Q1A's single-axis test left the residual $\chi^2 \approx 952$
essentially unmoved. Q1B is the registered joint multi-axis test
of whether stratifying further on a second $Z$ atop $Z_c$ closes
the gap. The experiment ran 2026-05-27; this section reports the
result.

\paragraph{Q1B: candidate set (head-to-head second $Z$ on top of
$Z_c$).}
With the cumulative baseline = Q1A's M3 with $\sigma$-algebra
$(\mathrm{day\_type}, Z_c) = 952.5167$ (NEAR-TIE caveat carryforward):
\begin{itemize}[leftmargin=2em]
\item $Z_{2,b}$: binary overnight indicator
$(\mathbb{1}[h \notin \{1, \ldots, 6\}])$. Captures the
overnight-vs-rest split with no within-day granularity.
\item $Z_{2,c}$: $4$-level categorical per P1's production
\verb|_time_of_day_label| at
\path{experiment/p1_patra_sen/__main__.py:81-96}: overnight
$h \in \{1, \ldots, 6\}$, morning ramp $h \in \{7, \ldots, 11\}$,
afternoon $h \in \{12, \ldots, 17\}$, evening $h \in \{18, \ldots, 24\}$.
\end{itemize}
Both refit M3 per
$(\mathrm{day\_type}, Z_c, Z_2)$ cell, $32$ cells for $Z_{2,c}$
and $16$ for $Z_{2,b}$. SMC particle propagation per cell;
pool-to-$(\mathrm{day\_type}, Z_c)$-parent fallback for cells with
fewer than $200$ pre-cutoff library rows (\textbf{not} pool to
$\mathrm{day\_type}$-only; pooling further would silently undo
Q1A's $Z_c$ carryforward and contaminate the cumulative $\chi^2$).
$32$ pre-cutoff cells with $n \ge 2364$; pool-fallback fires zero
events. \emph{Two} BLOCKING reproduction-check assertions: M3
no-$Z$ vs Q2A $953.84$ at tol $50$; M3 $\sigma = (\mathrm{day\_type},
Z_c)$ vs Q1A $952.5167$ at tol $50$ --- the cumulative-baseline
anchor.

\paragraph{Q1B: pre-registered outcomes.}
Same cuts as Q1A (inherited from the thread skeleton).
$\mathrm{better}\_\chi^2 = \min(\chi^2_{Z_{2,b}}, \chi^2_{Z_{2,c}})$
on post-cutoff marginal PIT, with R-A1B $\le 228$ / R-B1B $\in
(228, 2284]$ / R-C1B $> 2284$. Substantive secondary:
$\Delta\chi^2_{\mathrm{cum}} = 952.5167 - \mathrm{better}\_\chi^2$
measures cumulative reduction above Q1A's $Z_c$ carryforward.

\paragraph{Q1B: result.}
With $1000$-paired-day bootstrap CIs:

\begin{center}
\begin{tabular}{lrrr}
\toprule
& point & CI low & CI high \\
\midrule
$\chi^2_{Z_{2,b}}$ (binary)         & $696.0$ & $492.3$ & $1{,}009.0$ \\
$\chi^2_{Z_{2,c}}$ ($4$-level)      & $662.2$ & $479.8$ & $945.4$ \\
$\mathrm{better}\_\chi^2$           & $662.2$ & $468.1$ & $944.5$ \\
$\chi^2_{\mathrm{diff}}$            & $+33.8$ & $-53.2$ & $+124.5$ \\
$\Delta\chi^2_{\mathrm{cum}}$       & $+290.3$ & $+8.06$ & $+484.4$ \\
\bottomrule
\end{tabular}
\end{center}

\textbf{Both reproduction checks EXACT}: M3 no-$Z$ $\to$ Q2A
$|\mathrm{diff}| = 0.003$; M3 $Z_c$-only $\to$ Q1A $|\mathrm{diff}| =
0.0000$. The strongest dual-anchor reproduction in the workflow's
history. Synthetic gate PASS including the false-$Z_2$ companion
(collapse to the $Z_c$-only marginal mixture under randomised
$Z_2$ labels).

\paragraph{Q1B: verdict (with explicit CI caveat).}
Mechanical: \textbf{R-B1B} ($\mathrm{better}\_\chi^2 = 662.2 \in
(228, 2284]$), winner $Z_{2,c}$ at NEAR-TIE strength
($|\chi^2_{\mathrm{diff}}|/\mathrm{SE} < 1$). The substantive
secondary $\Delta\chi^2_{\mathrm{cum}} = +290.3$ at the point
estimate is \emph{$40\%$ of the remaining gap}, equivalent to
$218\times$ the closure Q1A delivered single-axis. \textbf{Substantive
verdict: R-B1B substantive FAVORABLE} on category, but with a
load-bearing caveat: the bootstrap CI $[+8.06, +484.4]$ on
$\Delta\chi^2_{\mathrm{cum}}$ spans a factor of roughly $60\times$.
The lower bound $+8.06$ is barely above Q1A's $+1.33$; the upper
bound $+484.4$ would close $67\%$ of the gap. The point estimate
is genuinely informative (it is positive and clear of zero at
$95\%$), but its magnitude is not. The framework reading
\emph{sharpens}, not refutes, S9+S10+S11+S12: single-axis bounded
near zero; \emph{joint} multi-axis bounded above zero with wide CI;
full closure to R-A1B does not fire even at upper CI ($662 - 484 =
178$ vs cut $228$ at upper-bound point projection, with the bootstrap
distribution itself wide). \emph{The identifiability obstruction
at full-closure level is preserved.}

\paragraph{(S13) The Q1B findings, three sub-claims.}
\begin{itemize}[leftmargin=2em]
\item \textbf{(S13a) Joint multi-axis within-data conditioning
recovers signal that single-axis could not, with wide CI.} The
$+290.3$ cumulative reduction at point exceeds Q1A's $+1.33$ by
$218\times$; lower-CI $+8.06$ still exceeds Q1A's reading. Joint
stratification is doing something single stratification was not.
The framework picture is sharpened: the $\sigma$-algebra
$\sigma(Z_c)$ left structure unabsorbed that
$\sigma(Z_c) \vee \sigma(Z_{2,c})$ partially captures. \emph{Magnitude
is genuinely uncertain across $60\times$.}

\item \textbf{(S13b) Per-cell decomposition exactly corroborates
P1, and the splits are heterogeneous within the P1 hot stratum.}
The saturday $Z_c = 6$ cell (Q1A $\chi^2 = 1634$; P1
$\hat\alpha_L = 0.61$) splits across $Z_{2,c}$ as
morning\_ramp $967 >$ evening $400 >$ afternoon $328$. The within-
P1-cell heterogeneity is large; the joint $\sigma$-algebra picks
up structure that \texttt{hour\_of\_week} alone could not resolve.

\item \textbf{(S13c) Methodologically: axis-level $\hat\alpha_L$
does NOT cleanly predict cell-level $\chi^2$ contribution under
joint stratification.} P1's axis-level $\hat\alpha_L = 0.11$ for
overnight predicted ``essentially Gaussian''; yet overnight
$(Z_{2,c} = 0)$ cells contribute substantially per cell
(sunday $Z_c = 6$ overnight $= 267$, saturday $Z_c = 5$
overnight $= 259$, weekday $Z_c \in \{1, \ldots, 4\}$ overnight
$= 174$--$211$). The axis $\to$ cell $\to$ $\chi^2$ mapping is
non-trivial under joint stratification. The proponent's ``P1
gradient explains per-cell ordering'' story failed structurally,
not just on magnitude. (S13c) is a methodological caveat for any
future per-cell $\chi^2$ interpretation derived from per-axis
$\hat\alpha_L$.
\end{itemize}

\paragraph{Forecaster calibration.}
Proponent (R-A1B, $\mathrm{better}\_\chi^2 = 180$, CI $[90, 450]$)
decisively falsified (realised $662$ is $1.47\times$ above CI
upper). Devil's advocate (R-B1B, $\mathrm{better}\_\chi^2 = 948$,
CI $[620, 1340]$) confirmed on outcome category but the realised
$662$ falls in the gap between proponent's CI upper ($450$) and
DA's CI lower ($620$). \textbf{Mixed scoring decomposition} ---
outcome and shape (NEAR-TIE preserved) from DA, direction
($Z_{2,c}$ wins) from proponent, magnitude from neither. Sixth
named-failure-mode instance, new shape variant.

\paragraph{Thread state after Q1B.}
Q1B.status $\to$ \texttt{settled} R-B1B; the thread's branching
rule for R-B1B is \texttt{null}, so mechanical state advancement
is unavailable. The advisory EXHAUSTED-UNRESOLVED pre-flag logged
at Q1A's settlement is \emph{downgraded}: Q1B's point-estimate
$+290.3$ refutes the pre-flag's premise (``single-axis $0.18\%$ at
binding axis $\Rightarrow$ multi-axis at weaker axis also $\sim 0\%$'').
The actual finding --- joint multi-axis recovers signal with wide
CI --- changes the next-decision shape: a three-way amendment-
direction decision is surfaced for user adjudication, with no
direction pre-committed by the coordinator: (a)~third within-data
$Z$ axis (\texttt{season} or \texttt{demand\_quantile}; weaker
prior than $Z_2$ was), (b)~new thread targeting cross-axis joint
stratification systematically, or (c)~temperature/exogenous-$Z$
arm (Q1C reopen or new thread per the framework-prescribed
resolution paths). Thread state \texttt{active}, current\_node =
\texttt{null} (paused). The right (a)/(b)/(c) call depends on
whether tightening Q1B's wide CI matters more than exploring the
next axis --- an advisor question. The CI lower bound $+8.06$ is
consistent with (a)/(b)/(c) all being roughly equipriotised; the
upper bound $+484.4$ would strongly favour (b) (the joint axis is
doing real work); neither extreme has prior empirical support
within the registered thread.

\section{Q1B' of the resolution-paths thread: settled (R-B1B', INSPECTION-ONLY, CI widened)}
\label{sec:q1b-prime-settled}

Q1B left a wide CI on the substantive secondary statistic
$\Delta\chi^2_{\mathrm{cum}}$. The user locked decision~(d) at
the end of session 2026-05-29-a: bracket-check the CI's source
before adjudicating the three-way amendment. The (a)/(b)/(c)
value depends on which CI endpoint is closer to truth, so tighten
first. Q1B' is the registered tightening run that asks whether
Q1B's wide CI is methodological (SMC particle-count sampling noise
at $N_{\mathrm{particles}} = 200$) or substantive (eval-window noise
on the $\sim 12{,}000$-row post-cutoff window plus paired-day
bootstrap finite-sample noise on a $\chi^2$ statistic of this
construction). Thread amendment A2 added Q1B' as a non-gating
sibling of Q1B under \verb|parent_branch = R-A| (mirroring the
A1 Q1A' pattern). The experiment ran 2026-05-29; this section
reports the result.

\paragraph{Q1B$'$: candidate set (single arm, parameter-only
variation).}
A single arm: re-run Q1B's winning $(\mathrm{day\_type}, Z_c,
Z_{2,c})$ configuration with $N_{\mathrm{particles}} = 400$ (vs
Q1B's $200$). All other methodology held bit-identical: same
pre-cutoff library, same \verb|mixture_2_gaussian_mle_fit|, same
SMC seed convention, same paired-day bootstrap convention (same
\verb|_Z2_BOOT_OFFSET = 8013|, $1000$ resamples), same
\verb|MIN_CELL_SIZE = 200| pool-fallback rule, same dual
reproduction-check anchors at Q2A $953.84$ and Q1A $952.5167$.
\textbf{Not} a head-to-head; \textbf{not} a methodology change;
the only knob is the SMC particle count.

\paragraph{Q1B$'$: pre-registered outcomes (verdict on CI WIDTH).}
The verdict statistic is the realised $95\%$ CI WIDTH on
$\Delta\chi^2_{\mathrm{cum}}$, with Q1B's settled width $476.358$
$\chi^2$ as the comparison baseline:
\begin{itemize}[leftmargin=2em]
\item \textbf{R-A1B$'$:} CI width $\le 357.27$ $\chi^2$
($\le 75\%$ of $476.358$), AND $|\mathrm{better}\_\chi^2_{\mathrm{point}}
- 662.223| \le 50$ $\chi^2$. SMC noise dominates the variance
(corresponds to $V_{\mathrm{smc}}/V_{\mathrm{total}} \ge 0.875$
under the bias-variance decomposition with
$V_{\mathrm{smc}} \to V_{\mathrm{smc}}/2$ at $N: 200 \to 400$);
the tightening lever works, user can adjudicate (a)/(b)/(c) on
the narrower CI.
\item \textbf{R-B1B$'$:} CI width $\in (357.27, 524.0]$ AND
point-shift $\le 50$. Mixed; SMC contributes but eval-window-
noise is non-trivial; the wide CI is partly intrinsic.
\item \textbf{R-C1B$'$:} CI width $> 524.0$ OR
$|\mathrm{point\text{-}shift}| > 50$. Widens or anomaly; SMC seed
dependence or methodological issue worth diagnosing before any
further amendment.
\end{itemize}

\paragraph{Q1B$'$: result.}
With $1000$-paired-day bootstrap CIs:

\begin{center}
\begin{tabular}{lrrr}
\toprule
& point & CI low & CI high \\
\midrule
$\chi^2_{Z_{2,b}}$                              & $703.235$ & $504.4$ & $1{,}019.1$ \\
$\chi^2_{Z_{2,c}}$                              & $679.128$ & $484.3$ & $991.9$ \\
$\mathrm{better}\_\chi^2$                       & $679.128$ & $480.9$ & $987.1$ \\
$\chi^2_{\mathrm{diff}}$ (paired SE $38.81$)    & $+24.107$ & $-48.2$ & $+103.7$ \\
$\Delta\chi^2_{\mathrm{cum}}$                   & $+273.388$ & $-34.576$ & $+471.622$ \\
\midrule
\textbf{CI WIDTH on $\Delta\chi^2_{\mathrm{cum}}$} & \multicolumn{3}{c}{$\mathbf{506.198}$} \\
\textbf{Q1B baseline width}                        & \multicolumn{3}{c}{$476.358$} \\
\textbf{Width change}                              & \multicolumn{3}{c}{$+29.84\ (+6.26\%)$} \\
$|\mathrm{better}\_\chi^2_{\mathrm{point}} - 662.223|$ & \multicolumn{3}{c}{$16.905$} \\
\bottomrule
\end{tabular}
\end{center}

Reproduction checks both within tol $50$: M3 no-$Z$ vs Q2A,
$|\mathrm{diff}| = 5.225$; M3 $Z_c$-only vs Q1A, $|\mathrm{diff}|
= 44.595$. Asymmetry ratio $8.54\times$.

\paragraph{Q1B$'$: verdict.}
Mechanical: \textbf{R-B1B$'$} (width $506.198 \in (357.27, 524.0]$
AND point-shift $16.905 \le 50$). Headroom above R-A1B$'$ cut:
$148.93$ $\chi^2$ (the corroboration is far from cleared); headroom
below R-C1B$'$ cut: $17.80$ $\chi^2$ (the falsification narrowly
avoided). The realised width sits at $89.3\%$ of the way through
the R-B1B$'$ band, $8.4\times$ closer to R-C1B$'$ than to R-A1B$'$;
a structurally different R-B1B$'$ shape than Q1B's settled R-B1B
which sat $\sim 2\times$ clear of both bounds.

Substantive verdict: \textbf{the CI WIDENED, did not narrow.}
Q1B$'$'s mechanical R-B1B$'$ category absorbs the realisation
because the band is wide enough to absorb both forecasters'
narrowing-direction forecasts and the realised widening; the band
is the unfalsifiability zone for direction errors. The substantive
reading is sharper than the mechanical category: SMC particle
doubling does \emph{not} tighten Q1B's CI, so the CI is
\emph{intrinsic to the data}, not to the SMC estimator. Q1B's wide
$[+8.06, +484.4]$ on $\Delta\chi^2_{\mathrm{cum}}$ reflects
eval-window noise on a $\chi^2$ statistic of this construction, not
SMC sampling noise.

\paragraph{(S14) The Q1B$'$ findings, three sub-claims.}
\begin{itemize}[leftmargin=2em]
\item \textbf{(S14a) SMC particle doubling does NOT tighten Q1B's
CI.} Width $476.358 \to 506.198$ (+$6.3\%$); point-shift
$16.9$ within tol. The proponent's strong-form SMC-dominance
forecast (width $338$ $[290, 395]$ requiring $V_{\mathrm{smc}}/
V_{\mathrm{total}} \ge 0.875$) is falsified $111$ $\chi^2$ above
its CI upper. The proponent's \emph{own} pre-registered
falsification condition (\#1: ``width $> 400 \Rightarrow$ strong-
form SMC-dominance refuted'') was triggered by the realisation.
The DA's modest-narrowing forecast (width $419$ $[370, 480]$,
requiring $V_{\mathrm{smc}}/V_{\mathrm{total}} \sim 0.07$) is
also falsified $26$ $\chi^2$ above its CI upper. Realised direction
implies $V_{\mathrm{smc}}/V_{\mathrm{total}} \ll 0.07$:
\textbf{eval-window noise dominates} Q1B's wide CI.

\item \textbf{(S14b) The reproduction-check side-channel
corroborates per-stratum SMC variance but refutes pooled-layer
SMC dominance.} No-$Z$ pooled drift $|\mathrm{diff}| = 5.2$; $Z_c$-
stratified parent drift $|\mathrm{diff}| = 44.6$ (ratio $8.54\times$).
SMC variance IS non-trivial at the $Z_c$-stratified parent layer
(doubling $N$ revealed it) but NOT at the no-$Z$ pooled layer.
This is a finite-sample per-stratum effect: $Z_c$ partitions
reduce per-cell $n$ from the full library to $\sim 1000$--$7000$
rows where SMC noise bites harder. DA's modest-SMC-variance
qualitative reading is corroborated; proponent's strong-form
SMC-dominance is refuted (if SMC truly dominated the CI, no-$Z$
drift would also be material at $N = 400$). (S14b) is also a
methodological caveat for (S10)'s per-$(h, \mathrm{day\_type})$
ablation analysis: the per-stratum SMC contribution that did not
surface at the $\sigma = (h, \mathrm{day\_type})$ granularity may
have biased (S10)'s share estimates inward.

\item \textbf{(S14c) Within-data tightening is empirically
impossible on this dataset.} Joint S9+S10+S11+S12+S13+S14 now
spans \emph{three independent layers} of identifiability-
obstruction evidence on the same residual $\chi^2$:
distributional-class (S9+S10 bounded), $\sigma$-algebra-enrichment
(S11+S12+S13 bounded with wide CI), methodological-tightening
(S14 SMC particle count cannot narrow the realised CI).
\textbf{Resolution requires auxiliary information beyond the
time index.} This is now a publication-grade conclusion, not a
hypothesis. The three-way (a)/(b)/(c) decision Q1B left for user
adjudication takes a S14-weighted prior: (c) temperature/exogenous-
$Z$ arm is the only direction not empirically refuted within the
registered thread.
\end{itemize}

\paragraph{Forecaster calibration.}
The seventh named-failure-mode instance, in a new shape variant:
both forecasters directionally falsified on the verdict-metric
central tendency. Mechanical R-B1B$'$ band absorbs the disagreement
because the band itself is wide enough; the registered band is
the unfalsifiability zone for direction errors. Proponent fails
outcome and falsifies its own pre-registered condition $\#1$; DA
passes outcome category but fails width point and direction.
Point-shift forecasts both succeed (realised $16.905$ inside
proponent $[2, 38]$ AND DA $[5, 47]$); DA's bias-correction
prediction of $\Delta\chi^2_{\mathrm{cum,point}} = +268$ (residual
$+5.4$) was the tightest single forecast surface. Distinct from
Q1B's mixed-scoring shape (direction-matched-one-side, magnitude-
matched-neither); Q1B$'$ has direction-matched-\emph{neither}.

\paragraph{Provenance grade.}
PROVENANCE-GRADE \textbf{INSPECTION-ONLY} per result.yaml
\verb|grade_demotion_reason|: the script's hardcoded
\path{experiment/results/q1b_z2_conditioning.txt} output path
collided with Q1B's CLAIM-GRADE artifact (Check P's substitution
audit set covered phase\_a hash + thread hash + prereg-dir but
missed the literal \texttt{.txt} filename at script line $1684$).
The run overwrote Q1B's committed CLAIM-GRADE artifact in-place;
restored from \texttt{HEAD} per user adjudication; demoted Q1B$'$
to INSPECTION-ONLY rather than launch a third script-edit + re-
audit cycle. The result.yaml and pickle ARE the durable record
(production-path code with all BLOCKING gates passed); downstream
citation must read ``inspection-only diagnostic.'' This is a
registry-housekeeping caveat, not a scientific one. A workflow-
design improvement is logged for the next-session retro: future
inline-edit parameterisations should grep-verify the result-
artifact \texttt{.txt} path string is in the substitution set
(suggested L1 item for Check P).

\paragraph{Thread state after Q1B$'$.}
Q1B$'$.status $\to$ \texttt{settled} R-B1B$'$; the thread's
\verb|branching_rules.Q1B'| has all three outcomes mapped to
\texttt{null} (non-gating per amendment A2), so mechanical state
advancement is unavailable and current\_node remains \texttt{null}
(unchanged from the S4 advance). Thread state \texttt{active};
the three-way (a)/(b)/(c) adjudication originally surfaced at
Q1B's S4 advance now proceeds on substantive grounds, with
(c)~temperature/exogenous-$Z$ arm priority upgraded by (S14c)
(only direction not empirically refuted within the registered
thread).

\section{What is settled at each level of generality}
\label{sec:settled-levels}

To summarise the campaign's standing under the kernel-family-aware
framework:

\begin{longtable}{P{0.40\textwidth}P{0.55\textwidth}}
\toprule
\textbf{Level} & \textbf{Status} \\
\midrule
\endfirsthead
\toprule
\textbf{Level} & \textbf{Status} (continued) \\
\midrule
\endhead
\midrule
\multicolumn{2}{r}{\textit{continued on next page}} \\
\endfoot
\bottomrule
\endlastfoot
Algebraic framework
(Defs.\ \ref{def:factorization}, \ref{def:coherence})
& \textbf{Settled} as a coherent formal proposal. Not yet
theory-side-audited for novelty against existing kernel-semigroup
literature. \\
\addlinespace
Working hypothesis on relaxed factorization being natural for
finite-dim embeddings of higher-dim processes
& \textbf{Empirically corroborated} in both Gaussian (S1, S2, S3)
and Student-$t$ (S8c) special cases. The defect persists under the
kernel-family lift, though at $\sim 10\times$ smaller magnitude than
the Gaussian estimate suggested. The hypothesis survives independent
of family. Bounded by (S9) and (S10) under both directions tested:
(S9) at fixed $Q$ via kernel refinement (M0 $\to$ M1 = $16\times$;
M1 $\to$ richer-best = $1.44\times$), (S10) via three pre-registered
same-$Q$ non-distributional mechanisms (all three below the $0.15$
null floor; two sign-determined negative). Further reduction of the
marginal pathology requires neither additional distributional
flexibility (S9) nor naive same-$Q$ post-hoc fixes (S10), but
$\sigma$-algebra enrichment via auxiliary information (see the
``how far the lift can go'' row below). \\
\addlinespace
Gaussian kernel family as the right family for this empirical
process
& \textbf{Empirically falsified at the marginal level} (S5).
Confirmed by (S8) under Q1: a Student-$t$ kernel reduces marginal
$\chi^{2}$ by $16\times$, but Student-$t$ alone does not fully
calibrate either. The conditional residual is heavy-tailed
($\hat\nu \approx 4.5$); the seasonal climatology is approximately
Gaussian (S8d). (S9) tightens this: even the most flexible
distributional family tested (KDE residual law) does not fully
calibrate either; the marginal pathology is not purely distributional. \\
\addlinespace
Specific structural prediction: defect concentrates in one
eigenmode with characteristic timescale $\tau$ (seed \S5.2)
& \textbf{Empirically falsified} (S2 shows the defect is diffuse
across eigenmodes). The framework permits structured defects;
this particular structure was wrong. \\
\addlinespace
Specific directional prediction: iterated under-estimates direct
diffusion at large $h$ (seed \S5.1)
& \textbf{Empirically falsified} (S6: iteration over-estimates;
saturation argument). Framework predicts a defect; the direction
was wrong. \\
\addlinespace
Specific magnitude prediction: $\mathrm{D\_RATIO} \approx 29$
as a stable headline of v2 (S1)
& \textbf{Empirically retained but qualified} by (S8c). The
$29\times$ figure is kernel-family-conditional; under the
corrected Student-$t$ kernel the analogous statistic is
$2.93$. v2's settled finding stands; the magnitude requires the
asterisk that $\sim 90\%$ was absorbing distributional-class
misspecification. \\
\addlinespace
Where heavy-tailedness lives in the process
(conditional residual vs.\ seasonal climatology)
& \textbf{Newly settled} by (S8d). The heavy tails live in the
$\hat\kappa_1$ innovation conditional on recent state, not in the
marginal demand given calendar. M1 outperforms M2; refining a
Student-$t$ climatology adds noise without capturing real signal. \\
\addlinespace
How far the distributional-class lift can go on this empirical
process
& \textbf{Empirically bounded} by (S9) on Lever~1 (kernel
refinement at fixed $Q$): three richer-than-Student-$t$ families
$\chi^{2}$-indistinguishable, KDE marginally worst; M0 $\to$ M1
$= 16\times$, M1 $\to$ Q2A-best $= 1.44\times$. \textbf{Further
bounded} by (S10) on the pre-registered Lever-$2$-like candidates:
mean-bias absorption, one-sided variance cap, and seam exclusion
are post-hoc same-$Q$ fixes (not true $\sigma$-algebra
enrichments); all three below the $0.15$ null floor with two
sign-determined negative (variance cap and seam exclusion
actively degraded calibration). Picture~(C) survives (the defect
remains; $\mathrm{D\_RATIO}_{\Studt} > 1$ with CI clear of $1$
per (S8c)). The joint S9~$+$~S10 pattern --- kernel refinement
plateaus AND naive same-$Q$ manipulations fail --- is applied
evidence for the identifiability obstruction (Bergna et al.\
$2026$ Prop.\ $1$; Heckman--Singer $1984$; Allahverdyan $2020$):
intrinsic heavy-tailedness and latent-$Z$ mixture artefacts are
observationally indistinguishable from $\kappa_{Q}$ alone. The
distributional-class line of inquiry is exhausted on its
pre-registered scope; the next move requires auxiliary information
per the framework's three resolution paths (candidate-$Z$ testing,
stratified NPMLEmix per Deb et al.\ 2022, synthetic Gaussian-DGP
calibration). \\
\addlinespace
Decomposition of the residual $\chi^{2}$ into the pre-registered
non-distributional mechanisms (iterated mean trajectory bias,
iterated variance divergence, day-anchor seam)
& \textbf{Empirically inadequate} by (S10). All three shares below
the $0.15$ null floor (point estimates $+0.109, -0.114, -0.139$);
two of three sign-determined NEGATIVE at $95\%$ CI (\texttt{variance\_div}
$\in [-0.22, -0.01]$, \texttt{seam} $\in [-0.17, -0.09]$). The
mech~$2$ variance cap and mech~$3$ seam exclusion actively WORSENED
$\chi^{2}$. Mech~$1$ mean-bias contributes $+0.11$ share but well
below the $0.50$ dominance threshold. Together with (S9), this
maps the boundary of the distributional-class framework's two
natural levers (kernel refinement at fixed $Q$; naive same-$Q$
manipulations of the predictive distribution). The remaining
$\sim 954$ $\chi^{2}$ requires $\sigma$-algebra enrichment with
auxiliary $Z$ (per the disintegration-as-diagnostic note parked
2026-05-26 in the theory-side programme:
\path{Resolvent_Framework/notes/unsorted/disintegration_diagnostic.md}).
The under-dispersion signature ($\sigma^{2}_{\mathrm{iter}} <
\sigma^{2}_{\mathrm{empirical}}$ at virtually every cell)
independently points at missing exogenous content. \\
\addlinespace
Where in the existing time-indexed $\sigma$-algebra the missing-$Z$
structure is concentrated (per-stratum Patra--Sen diagnostic)
& \textbf{Empirically localised} by (S11). The per-stratum
$\hat\alpha_L$ diagnostic on Q1's M1 PIT residuals fires R-A at
landslide strength: $\max_{\mathrm{axis}} \mathrm{range}(\hat\alpha_L)
= 0.575$ $[0.4475, 0.6475]$, lower CI clearing the $0.20$ cut by
more than $2\times$. The binding axis is \texttt{hour\_of\_week},
not \texttt{season} (which ranks fourth of five at $0.135$). The
missing-$Z$ structure is \emph{cyclical/social-weekly, not
seasonal/climatological}: hot zone is the weekend window (Fri
10:00 -- Sun 23:00) at $\hat\alpha_L \in \{0.55, 0.61, 0.55\}$;
cold zone is Monday morning at $\hat\alpha_L = 0.035$ (essentially
Gaussian). The temperature prior from
\path{mitacs-clim-gap-nonstationary} takes a substantive downward
update on its localising-axis claim. Day-type ranking (saturday
$0.6475 >$ sunday $0.4850 >$ weekday $0.3025$) corroborates Q2A
Finding~$2$ and eigenmodes-v1 (S2, S3) from a distinct estimator
family. \\
\addlinespace
How far $\sigma$-algebra enrichment on the localised axis can
go on this empirical process
& \textbf{Empirically bounded near zero at single-axis} by (S12),
\textbf{bounded above zero with wide CI at joint multi-axis} by
(S13), and \textbf{empirically un-tightenable from the SMC
estimator side} by (S14). (S12) re-fits M3 with the binding axis
(\texttt{hour\_of\_week}) $Z$ added at $\sigma = (\mathrm{day\_type},
Z_c)$ and shaves $1.33$ $\chi^2$ off Q2A's baseline $953.84$: $0.18\%$
of the $725.84$ gap-to-close. The contamination is real and
localised where P1 said it would be (saturday $\mathrm{bin}_6$
contributes $\chi^2 = 1634$ pre-stratification, exactly the
$\hat\alpha_L = 0.61$ stratum), but the M3 mixture-$2$ Gaussian
family is at capacity within strata --- combined with (S9)'s
$\chi^2$-indistinguishability of M3/M4/M5, no within-$Z$ same-class
substitution closes this cell. (S13) re-fits at $\sigma =
(\mathrm{day\_type}, Z_c, Z_{2,c})$ with $Z_{2,c}$ from P1's
second-strongest axis (\texttt{time\_of\_day}): point reduction
$+290.3$ $\chi^2$ ($40\%$ of remaining gap) but CI $[+8.06, +484.4]$
spans a factor of roughly $60\times$; full closure to R-A1B does
not fire even at upper CI. (S14) brackets the CI's source by
doubling SMC particle count from $N = 200$ to $N = 400$: the CI
width is $506.198$ vs Q1B's $476.358$ ($+6.3\%$, slightly wider),
proponent's strong-form SMC-dominance forecast is falsified by
its own pre-registered condition, eval-window noise dominates
the realised CI. Joint S9+S10+S11+S12+S13+S14 = three independent
layers of identifiability-obstruction evidence (distributional-
class, $\sigma$-algebra-enrichment, methodological-tightening).
\textbf{Resolution requires auxiliary information beyond the time
index}; this is now a publication-grade conclusion. The within-
data lever space is empirically exhausted on this dataset; the
next move is the (a)/(b)/(c) adjudication described in
\S\ref{sec:q1b-prime-settled}'s thread-state paragraph, with
S14-weighted prior on (c) (temperature/exogenous-$Z$ arm). \\
\addlinespace
Methodological caveat: axis-level $\hat\alpha_L$ as a predictor
of cell-level $\chi^2$ contribution under joint stratification
& \textbf{Newly settled (negative) by (S13c)}. Per-cell
decomposition of Q1B shows overnight ($Z_{2,c} = 0$) cells
contribute substantially per cell ($174$--$267$ $\chi^2$ per cell
across day-types) despite P1's axis-level $\hat\alpha_L = 0.11$
for overnight predicting ``essentially Gaussian.'' The axis
$\to$ cell $\to$ $\chi^2$ mapping is non-trivial under joint
stratification. Any future per-cell $\chi^2$ interpretation
derived from per-axis $\hat\alpha_L$ requires explicit
within-stratum verification. \\
\addlinespace
Per-stratum SMC variance contribution to per-cell $\chi^2$
estimates
& \textbf{Newly surfaced by (S14b)}. The reproduction-check
asymmetry at $N = 400$ (no-$Z$ pooled drift $5.2$ vs $Z_c$-
stratified parent drift $44.6$; ratio $8.54\times$) shows SMC
variance is non-trivial at the per-stratum layer where per-cell
$n$ is $1000$--$7000$ rows, but not at the no-$Z$ pooled layer.
This is a methodological caveat for (S10)'s per-$(h,
\mathrm{day\_type})$ ablation share estimates: the per-stratum
SMC contribution that did not surface at that granularity may
have biased the share estimates inward. Not a sign-flip on
(S10) (mechanism shares are bounded by the per-stratum SMC
contribution only in magnitude); but a flag for future
within-stratum analyses to budget for per-cell SMC noise
explicitly. \\
\addlinespace
Test methodology (synthetic-calibrated divergence test) as a
generic discipline
& \textbf{Settled} across five cycles. Generalises to non-Gaussian
families with no structural change to the discipline; only the
divergence and family change. \\
\addlinespace
Workflow (pre-register, code-path-audit, phase-fidelity, arbiter, thread)
& \textbf{Settled} across two full threads plus four single-
experiment cycles and one non-gating sibling under amendment.
Catches the named retraction modes reliably; recognises
pre-registered-criterion-vs-legitimate-data-structure as the
fourth named failure mode (now \emph{seven} instances spanning
multiple shape variants: eigenmodes-v1 O2-vs-O3 near-tie,
PIT-calendar-fingerprint P1 tautology, Q1 R-D-vs-R-C MLE
saturation, Q2B mech\_2/mech\_3 sign-determined negative shares,
P1 mechanical-matches-category-with-within-bucket-detail-matching-
neither, Q1A DA-CONFIRMED-with-exact-within-bucket-point-but-
wrong-direction, Q1B mixed-scoring decomposition, Q1B$'$
both-forecasters-directionally-falsified-on-width-band-absorbs-
disagreement). The arbiter pattern IS the prevention; no
engineering fix attempted across the seven instances. \\
\addlinespace
Workflow lab-discipline cost shape: ratio of reactive cascades
to intended end-to-end steps
& \textbf{Newly characterised by Q1B$'$ session retro
(2026-05-29)}. Q1B$'$'s six-step end-to-end accumulated four
reactive cascades (cosmetic \texttt{>=} $\to$ \texttt{>}
amendment in $5$ re-stamps; script parameterisation $10$-string
substitution; gitignore for sibling pickle; Check P substitution-
audit miss on the result-artifact \texttt{.txt} filename leading
to in-place overwrite of Q1B's CLAIM-GRADE artifact and
user-adjudicated demote to INSPECTION-ONLY). Two patterns logged
for future workflow-rule refinement: (i) cosmetic-amendment-
cost-asymmetry --- the carve-out's per-cascade cost exceeds the
realised benefit when the underlying defect has probability-zero
realised impact (proposed rule: cosmetic typos with PASS\_WITH\_
ADVISORY-grade Check T outcomes may stay ADVISORY-only); (ii)
the script's artifact-pinning strings (prereg-dir, phase\_a
hash, thread hash, output-artifact \texttt{.txt} filename) need
a single source of truth (proposed refactor: \texttt{--prereg-dir}
CLI flag $\to$ derive paths/hashes mechanically). Neither has
been retracted as a result; both are workflow-design
optimisations whose absence accumulates session-time cost
linearly with sibling-experiment count. \\
\end{longtable}

The framework's general formulation in \S\ref{sec:framework} is
the right object to register with the theory-side audit pipeline
(currently the Resolvent\_Framework programme's
\verb|notes/unsorted/| inbox). Its novelty against the
SDE-generator-inference, Koopman--EDMD, and renormalization-group
literatures has not been audited; the seed flagged this gap and
left it to the theory-side workflow.

\appendix

\section{Provenance and references}
\label{app:provenance}

\paragraph{Registry entries.}
Per-experiment artifacts under
\verb|notes/preregistrations/<date>_<topic>/|. All entries
hash-stamped via \verb|experiment.audit.registry.stamp(...,| \verb|self_path=...)|;
chains verifiable via \verb|registry verify|.

\noindent
\begin{tabular}{P{0.12\textwidth}P{0.53\textwidth}P{0.27\textwidth}}
\toprule
\textbf{Source} & \textbf{Entry} & \textbf{Verdict} \\
\midrule
S1, S6
  & \path{2026-05-26_multiscale-direct-h-step-v2}
  & SETTLED \\
\addlinespace
S2, S3
  & \path{2026-05-26_multiscale-defect-eigenmodes}
  & SETTLED MIXED \\
\addlinespace
S4, S5
  & \path{2026-05-26_pit-calendar-fingerprint}
  & SETTLED AMBIGUOUS \\
\addlinespace
S8 (Q1)
  & \path{2026-05-26_q1-student-t-vs-gaussian}
  & SETTLED (R-D mechanical / R-C substantive) \\
\addlinespace
S9 (Q2A)
  & \path{2026-05-26_q2a-richer-family-mixture-or-nonparametric}
  & SETTLED R-B2 \\
\addlinespace
S10 (Q2B)
  & \path{2026-05-26_q2b-non-distributional-decomposition}
  & SETTLED R-C3 \\
\midrule
Thread \newline (Gaussian)
  & \texttt{2026-05-26\_missing-} \newline \texttt{content-thread}
  & \texttt{exhausted} \\
\addlinespace
Thread \newline (kernel-family)
  & \texttt{2026-05-26\_distributional-} \newline \texttt{class-thread}
  & \texttt{exhausted} per (S9)~$+$~(S10); fresh thread on the framework's three resolution paths is the natural next move \\
\bottomrule
\end{tabular}

\paragraph{Supporting memories.}
The lab memory store lives under
\path{~/.claude/projects/}; the project-specific subdirectory holds
each named memory as a markdown file. Memories referenced above:

\begin{itemize}[leftmargin=2em]
\item \texttt{mitacs-rebaseline-facts} --- production-validated
in-sample non-Gaussianity (excess kurtosis $\sim 24$--$33$); the
empirical basis for the $\nu \approx 4.5$ forecast in S8a.
\item \texttt{mitacs-rank1-structural} --- $\Sigma_1$ is rank-one
by construction of the single-variable delay embedding; the only
stochastic content is the coord-$0$ innovation.
\item \texttt{mitacs-theta-rail-pinning} --- the estimator-as-run
reduces to global OLS at $\theta = 0$; bandwidth localisation is
empirically non-discriminating.
\item \texttt{mitacs-clim-gap-nonstationary} --- the post-cutoff
demand drift that S4's secular-drift candidate corroborates.
\item \texttt{mitacs-multiscale-direct-h-step-v2} --- the SETTLED
v2 finding; its magnitude now qualified by S8c.
\item \texttt{mitacs-multiscale-defect-eigenmodes-v1} --- S2, S3.
\item \texttt{mitacs-pit-calendar-fingerprint-p1} --- S4, S5.
\item \texttt{mitacs-q1-student-t-vs-gaussian} --- S8.
\item \texttt{mitacs-q2a-richer-family} --- S9.
\item \texttt{mitacs-q2b-non-distributional} --- S10.
\item \texttt{mitacs-session-lessons-corpus} --- the named
retraction modes; the fourth mode (criterion-vs-legitimate-data-
structure) added with Q1 as the third instance; Q2B adds the
fourth instance with substantive verdict \emph{stronger} than
mechanical (the prior three had it weaker/different).
\end{itemize}

\paragraph{Seed and workflow.}
\begin{itemize}[leftmargin=2em]
\item This memo
(\path{writeup/tex/missing_content_memo.tex}) is itself the
kernel-family-agnostic statement of the framework that the
original Gaussian-implicit seed
(\path{notes/seeds/multiscale_factor_coherence.md}) prompted. The
seed remains in the registry for lineage; this document supersedes
its formal content.
\item Workflow design:
\path{notes/seeds/applied_audit_workflow.md}.
Agents (\path{.claude/agents/}):
\texttt{code-path-auditor},
\texttt{phase-fidelity},
\texttt{preregister},
\texttt{devils-advocate},
\texttt{arbiter},
\texttt{thread-coordinator},
plus
\texttt{multiverse},
\texttt{framework-escalator},
\texttt{pre-experiment-checklist},
\texttt{editorial-pass},
\texttt{literature-scout}.
Dispatcher: \path{.claude/skills/audit/SKILL.md}.
Registry implementation: \path{experiment/audit/registry.py}
(CLI \texttt{registry list}, \texttt{registry verify},
\texttt{registry thread status}, etc.).
\item Q1's experiment script:
\path{experiment/distributional_class_q1/__main__.py}.
Student-$t$ MLE in the production whitelist:
\path{processing.innovations.estimator.student_t_mle_fit}.
Output artifact:
\path{scratch/data/distributional_class_q1/q1.pkl}.
\item Q2A's experiment script:
\path{experiment/distributional_class_q2a/__main__.py}.
Richer-family fitters in the production whitelist:
\path{processing.innovations.estimator.mixture_2_gaussian_mle_fit},
\path{...mixture_3_gaussian_mle_fit},
\path{...kde_residual_fit}.
Q2A synthetic gate (METHOD/DESIGN):
\path{experiment/results/q2a_synthetic_gate.txt}.
Output artifact:
\path{scratch/data/distributional_class_q2a/q2a.pkl}.
\item Q2B's experiment script:
\path{experiment/distributional_class_q2b/__main__.py}.
Primary CLAIM artifact:
\path{experiment/results/q2b_non_distributional.txt}.
Output pickle (per-mechanism PIT tables, per-cell attribution-only
statistics):
\path{scratch/data/distributional_class_q2b/q2b.pkl}.
\item Theory-side framework note parked 2026-05-26 identifying the
identifiability obstruction and the three resolution paths:
\path{disintegration_diagnostic.md} in
\path{Resolvent_Framework/notes/unsorted/}.
\end{itemize}

\paragraph{Next concrete step.}
The \texttt{distributional-class-thread} was declared
\texttt{exhausted} on 2026-05-26 (commit \texttt{543e3e8})
following Q2B's R-C3 verdict; under the framework reframe
(S9~$+$~S10 jointly bound the two distributional-class levers),
amending the thread further within its pre-registered scope
would not address the identifiability obstruction. The natural
next move is a fresh thread on the framework's three resolution
paths: (i) candidate-$Z$ testing (temperature, secular drift
via \texttt{year\_since\_cutoff} per (S5), refined day-type
partition); (ii) stratified covariate-dependent two-component
mixture decomposition via NPMLEmix (Deb et al.\ 2022) as a
$\sigma$-algebra adequacy diagnostic; (iii) synthetic Gaussian-DGP
calibration to anchor the residual-$\chi^{2}$ floor.

\end{document}
